\chapter{存储层次与访问全景}

\begin{keyidea}
存储系统不是一排容量和速度数字，而是一条由不同粒度、协议和完成语义组成的访问路径。本章先建立从 load 到块 I/O 的全局坐标系。
\end{keyidea}

CPU 寄存器堆里一个 bit 由两个交叉耦合反相器保存，主存里一个 bit 由一颗电容保存，闪存里一个 bit 由被困在绝缘层中的电子保存，硬盘里一个 bit 由一小块磁性薄膜的磁化方向保存。同样是 0 和 1，物理载体完全不同，读出方式、写入代价、保持时间、访问粒度和失效模式也完全不同。本章把这四类存储器件的内部结构讲到可以画截面图、算噪声裕量、读时序参数的程度。

\begin{keyidea}
存储层次不是“快慢不同的同一种东西”，而是\hwkey{物理机制完全不同的器件}。SRAM 靠反馈锁存，DRAM 靠电容电荷，NAND 靠绝缘层中的被困电子，HDD 靠磁畴方向。每一层的访问粒度、写入代价和失效模式都由其物理结构决定。
\end{keyidea}

\section{一、从 bit 到字到阵列到系统：结构递进}

所有存储器件都沿同一条结构递进线扩展：先有一个能保存 1 bit 的物理点，再把多个点排成一行形成一个字（word），再把多行多列排成二维阵列（array），最后把多个阵列在第三维度上堆叠或并列组成系统（bank、plane、die、channel、platter）。每一层的扩展方式不同，但递进逻辑相同。

这里先固定两个贯穿全章的布线术语。\term{字线}{word line, WL}沿行方向走线，一条字线有效即选中该行所有单元；\term{位线}{bit line, BL}沿列方向走线，负责把选中单元的数据传送出去（读）或把外部数据送入单元（写）。把阵列排成二维而不是一长串，是译码代价使然：直接选择 $N$ 个存储点需要 $N$ 条选择线，排成 $\sqrt{N}\times\sqrt{N}$ 方阵后，行、列各 $\sqrt{N}$ 条线即可定位任意一点，选择线总数从 $N$ 降为 $2\sqrt{N}$，译码器和驱动器的规模随之指数级缩小。系统一级的名词将在对应器件的章节展开，此处先记住轮廓：bank 是可独立接收命令的子阵列组，die 是晶圆上切下的单颗芯片，channel 是可独立传输的数据通道，platter 是磁盘中的盘片。

\begin{pdfdiagram}[x=1cm,y=1cm]
  % === Level 1: 单点 (single cell) ===
  \node[hw label, font=\small\bfseries] at (1.2,3.8) {单点 (1 bit)};
  % Draw a small capacitor symbol representing one storage cell
  \draw[BookTeal, line width=0.7pt] (0.6,2.8) -- (0.6,2.2);
  \draw[BookTeal, line width=0.9pt] (0.3,2.2) -- (0.9,2.2);
  \draw[BookTeal, line width=0.9pt] (0.3,1.9) -- (0.9,1.9);
  \draw[BookTeal, line width=0.7pt] (0.6,1.9) -- (0.6,1.3);
  % Small transistor gate symbol
  \draw[BookAccent, line width=0.6pt] (1.2,2.4) -- (1.2,1.7);
  \draw[BookAccent, line width=0.6pt] (1.0,2.4) -- (1.4,2.4);
  \draw[BookAccent, line width=0.6pt] (1.0,1.7) -- (1.4,1.7);
  \draw[BookAccent, line width=0.6pt] (1.4,2.05) -- (1.8,2.05);
  % Labels
  \node[hw label, font=\tiny] at (0.6,1.0) {$C_S$};
  \node[hw label, font=\tiny] at (1.8,2.4) {WL};
  \node[hw label, font=\tiny, text width=2.2cm] at (1.2,0.5) {一个电容/锁存器/浮栅/磁畴};
  % Arrow to next level
  \draw[BookRule, line width=0.6pt, ->] (2.3,2.0) -- (2.9,2.0);

  % === Level 2: 线性 (one word) ===
  \node[hw label, font=\small\bfseries] at (4.5,3.8) {线性 (1 word)};
  % Shared wordline (horizontal)
  \draw[BookAmber, line width=0.8pt] (3.2,3.0) -- (5.8,3.0);
  \node[BookAmber, font=\tiny] at (3.0,3.0) {WL};
  % 4 cells in a row
  \foreach \x in {3.5, 4.1, 4.7, 5.3} {
    \draw[BookTeal, line width=0.5pt] (\x,3.0) -- (\x,2.4);
    \draw[BookTeal, line width=0.6pt] (\x-0.15,2.4) -- (\x+0.15,2.4);
    \draw[BookTeal, line width=0.6pt] (\x-0.15,2.2) -- (\x+0.15,2.2);
    \draw[BookTeal, line width=0.5pt] (\x,2.2) -- (\x,1.6);
  }
  % Bitlines (vertical, going down)
  \foreach \x in {3.5, 4.1, 4.7, 5.3} {
    \draw[BookAccent, line width=0.4pt] (\x,1.6) -- (\x,1.0);
  }
  \node[BookAccent, font=\tiny] at (5.6,1.0) {BL};
  % Dots for more bits
  \node[hw label, font=\tiny] at (5.9,2.3) {$\cdots$};
  \node[hw label, font=\tiny, text width=2.4cm] at (4.5,0.5) {8--64 bit 共享 WL，BL 并行读出};
  % Arrow to next level
  \draw[BookRule, line width=0.6pt, ->] (6.0,2.0) -- (6.6,2.0);

  % === Level 3: 平面 (2D array) ===
  \node[hw label, font=\small\bfseries] at (8.2,3.8) {平面 (2D 阵列)};
  % Grid of cells (4x4)
  \foreach \x in {7.2, 7.7, 8.2, 8.7} {
    \foreach \y in {1.2, 1.7, 2.2, 2.7} {
      \fill[BookTeal!40] (\x,\y) circle (0.08);
    }
  }
  % Wordlines (horizontal)
  \foreach \y in {1.2, 1.7, 2.2, 2.7} {
    \draw[BookAmber, line width=0.3pt] (7.0,\y) -- (9.0,\y);
  }
  % Bitlines (vertical)
  \foreach \x in {7.2, 7.7, 8.2, 8.7} {
    \draw[BookAccent, line width=0.3pt] (\x,1.0) -- (\x,2.9);
  }
  % Row decoder (left)
  \draw[fill=BookAmber!15, line width=0.4pt] (6.7,1.0) rectangle (7.0,2.9);
  \node[hw label, font=\tiny, rotate=90] at (6.55,2.0) {行译码};
  % Column mux (bottom)
  \draw[fill=BookAccent!15, line width=0.4pt] (7.0,0.7) rectangle (9.0,1.0);
  \node[hw label, font=\tiny] at (8.0,0.85) {列 mux};
  \node[hw label, font=\tiny, text width=2.4cm] at (8.2,0.2) {行$\times$列矩阵，二维寻址};
  % Arrow to next level
  \draw[BookRule, line width=0.6pt, ->] (9.3,2.0) -- (9.9,2.0);

  % === Level 4: 立体 (3D system) ===
  \node[hw label, font=\small\bfseries] at (11.5,3.8) {立体 (系统)};
  % 3D stacked blocks (isometric-ish)
  % Bottom layer
  \draw[fill=BookTeal!20, line width=0.5pt] (10.2,0.8) -- (11.4,0.8) -- (12.0,1.2) -- (10.8,1.2) -- cycle;
  \draw[fill=BookTeal!30, line width=0.5pt] (10.2,0.8) -- (10.2,1.1) -- (10.8,1.5) -- (10.8,1.2) -- cycle;
  \draw[fill=BookTeal!15, line width=0.5pt] (10.2,1.1) -- (11.4,1.1) -- (12.0,1.5) -- (10.8,1.5) -- cycle;
  % Middle layer
  \draw[fill=BookAccent!20, line width=0.5pt] (10.2,1.5) -- (11.4,1.5) -- (12.0,1.9) -- (10.8,1.9) -- cycle;
  \draw[fill=BookAccent!30, line width=0.5pt] (10.2,1.5) -- (10.2,1.8) -- (10.8,2.2) -- (10.8,1.9) -- cycle;
  \draw[fill=BookAccent!15, line width=0.5pt] (10.2,1.8) -- (11.4,1.8) -- (12.0,2.2) -- (10.8,2.2) -- cycle;
  % Top layer
  \draw[fill=BookNavy!20, line width=0.5pt] (10.2,2.2) -- (11.4,2.2) -- (12.0,2.6) -- (10.8,2.6) -- cycle;
  \draw[fill=BookNavy!30, line width=0.5pt] (10.2,2.2) -- (10.2,2.5) -- (10.8,2.9) -- (10.8,2.6) -- cycle;
  \draw[fill=BookNavy!15, line width=0.5pt] (10.2,2.5) -- (11.4,2.5) -- (12.0,2.9) -- (10.8,2.9) -- cycle;
  % Labels for layers
  \node[hw label, font=\tiny] at (12.4,1.0) {bank 0};
  \node[hw label, font=\tiny] at (12.4,1.7) {bank 1};
  \node[hw label, font=\tiny] at (12.4,2.4) {bank $n$};
  % Channel annotation
  \draw[BookRed, line width=0.4pt, <->] (10.0,0.8) -- (10.0,2.5);
  \node[BookRed, font=\tiny, rotate=90] at (9.7,1.6) {通道};
  \node[hw label, font=\tiny, text width=2.4cm] at (11.5,0.2) {bank/plane/die/通道/盘片};
\end{pdfdiagram}
\diagramnote{存储器件的四级结构递进。左起：单个存储元件（电容/锁存器）→多个元件共享字线排成一行→多行多列组成二维阵列（行译码+列 mux）→多个阵列在第三维度堆叠或并列组成系统。每一级引入新的寻址维度和新的物理约束。}

下表把四类器件在每一级的具体实现并列对照：

\begin{longtable}{L{0.10\linewidth}L{0.22\linewidth}L{0.22\linewidth}L{0.22\linewidth}L{0.14\linewidth}}
\toprule
层级 & SRAM & DRAM & NAND Flash & HDD \\
\midrule
单点 (1 bit) & 6T 交叉耦合锁存 & 1T1C 电容 & 浮栅/电荷俘获 MOSFET & 磁畴（$\sim$数百颗粒） \\\tablerowrule
线性 (1 word) & 8--64 bit 共享 WL，BL 对并行 & 一整行（8--16 Kbit）共享 WL & 32--128 单元串联成 string & 一条磁道上的连续扇区 \\\tablerowrule
平面 (2D array) & 行$\times$列矩阵，行译码选 WL，列 mux 选字 & bank 内行$\times$列，row buffer 缓存整行 & plane 内 page$\times$block 矩阵 & 盘片表面：磁道$\times$扇区 \\\tablerowrule
立体 (系统) & 多 bank、多 port、多通道 cache & 多 bank group、多 rank、多通道 DIMM & 多 plane、多 die、多通道、3D 垂直堆叠 & 多盘片、多磁头、多区（ZBR） \\
\bottomrule
\end{longtable}

理解这条递进线后，后续每节的结构都按同一顺序展开：先讲单点物理（这个 bit 怎么存、怎么读、怎么写），再讲线性组织（多个 bit 怎么排成一行、共享什么信号），再讲平面阵列（行和列怎么寻址、怎么选中目标），最后讲立体系统（多个阵列怎么并列、怎么调度、怎么对外呈现接口）。

\begin{keyidea}
整机硬件不是“CPU 加一块内存”的粗框图，而是由一组定义明确的事务组成。每次访问都必须说明地址、方向、字节选择、目标片选、握手、完成、错误和副作用。只有这些边界明确，CPU 才能可靠地取指、访存、访问设备和处理中断。
\end{keyidea}

\begin{pdfdiagram}
  \node[hw state, minimum width=1.55cm] (cpu) at (0,0) {CPU};
  \node[hw compute, minimum width=1.95cm] (dec) at (2.35,0) {地址译码};
  \node[hw state, minimum width=1.65cm] (rom) at (5.0,1.0) {ROM};
  \node[hw state, minimum width=1.65cm] (ram) at (5.0,0) {SRAM};
  \node[hw control, minimum width=1.65cm] (io) at (5.0,-1.0) {MMIO};
  \node[hw control, minimum width=1.75cm] (irq) at (7.45,-1.0) {中断/DMA};
  \draw[hw bus] (cpu) -- node[above,hw label]{地址/控制} (dec);
  \draw[hw bus] (dec.east) |- (rom.west) node[pos=0.25,above,hw label]{片选};
  \draw[hw bus] (dec) -- (ram);
  \draw[hw bus] (dec.east) |- (io.west);
  \draw[hw return] (io) -- (irq);
  \draw[hw return] (irq.south) -- ++(0,-0.5) -| (cpu.south) node[pos=0.35,below,hw label]{完成/事件};
  \node[hw label, text width=8.8cm] at (3.75,-2.45) {整机事务从 CPU 发起，但完成、等待和错误由目标与互连共同决定；设备还可通过中断或 DMA 把事件带回 CPU。};
\end{pdfdiagram}
\diagramnote{从最小 CPU 到整机事务。地址译码决定目标，握手和事件决定本次访问何时真正完成。}

\section{二、主存不是抽象数组}

软件把内存看成按地址编号的字节序列，硬件却必须用地址线、译码器、存储阵列、数据线和控制信号实现这件事。一次读访问至少包含四个事实：发起者给出地址，译码逻辑选中目标，目标经过内部访问后驱动数据，发起者在规定边界采样。一次写访问则要求地址、写数据、字节使能和写控制在目标采样窗口内稳定。

\begin{longtable}{L{0.18\linewidth}L{0.34\linewidth}L{0.38\linewidth}}
\toprule
对象 & 硬件含义 & 关键问题 \\
\midrule
地址 & 指出本次事务目标位置 & 高位译码是否只选中一个目标 \\\tablerowrule
数据 & 读时由目标返回，写时由发起者提供 & 同一时刻是否只有一个驱动者 \\\tablerowrule
读写控制 & 区分本次事务方向和类型 & 输出使能和写使能是否互斥 \\\tablerowrule
字节使能 & 指明哪些 8-bit 车道有效 & 字节写不会破坏相邻数据 \\\tablerowrule
ready/valid & 指明请求或响应是否真正完成 & 慢目标能否插入等待而不丢地址和数据 \\
\bottomrule
\end{longtable}

\topic{ROM、SRAM 和 DRAM 的读写不是同一种动作}
ROM 或 Flash 适合保存启动代码，掉电后内容仍在；SRAM 接口直观，地址稳定后在访问时间内返回数据，适合教学板、cache 和小容量片上 RAM；DRAM 密度高，但单元需要刷新，访问还要经过 bank、行、列、预充电和控制器调度。把这些对象都叫“内存”没有错，但在硬件设计中必须区分它们的时序和控制边界。

\begin{longtable}{L{0.18\linewidth}L{0.34\linewidth}L{0.38\linewidth}}
\toprule
存储类型 & 典型用途 & 关键边界 \\
\midrule
ROM/Flash & 复位入口、固件、常量表 & 复位地址必须映射到有效内容，读时序满足取指 \\\tablerowrule
SRAM & cache、实验板 RAM、小容量片上存储 & \code{CS}/\code{OE}/\code{WE} 互斥，地址和数据满足建立时间 \\\tablerowrule
DRAM & 大容量主存 & 控制器负责 activate、read/write、precharge、refresh \\\tablerowrule
寄存器文件 & CPU 内部近距离状态 & 多读端口、写端口、旁路和同周期读写规则明确 \\
\bottomrule
\end{longtable}

\topic{异步 SRAM 是最好的第一块主存}
教学整机可以先使用低速异步 SRAM。读周期中，地址和片选稳定后，打开输出使能，SRAM 在访问时间后驱动数据线；写周期中，输出使能关闭，发起者驱动数据线，在写使能有效窗口内把数据写入目标字节。这个模型足够简单，也足以暴露硬件最常见错误：片选重叠、输出争用、写使能过早、字节使能错误。

\begin{lstlisting}
// 注释（推进）：逐行追踪发起者、所有权和可见性；请求被接受不等于事务完成。
// 注释（边界）：以采样沿、状态位或 completion 为准，再允许消费者使用结果。
SRAM read:
  address stable
  CS active, OE active, WE inactive
  wait access time
  sample data

SRAM write:
  address and write_data stable
  CS active, OE inactive, WE active for write window
  selected byte lanes store data
\end{lstlisting}

\begin{pdfdiagram}[x=1cm,y=1cm]
  % ADDR — 总线由无效翻转到有效
  \draw[BookTeal,line width=0.85pt] (0.6,1.7) -- (1.6,1.2) -- (8.4,1.2);
  \draw[BookTeal,line width=0.85pt] (0.6,1.2) -- (1.6,1.7) -- (8.4,1.7);
  \node[hw label, anchor=east] at (0.5,1.45) {ADDR};
  \node[hw label] at (5.0,1.45) {地址稳定};
  % CS_n/OE_n — 低有效: 高→低→高
  \draw[BookAccent,line width=0.85pt] (0.6,0.55) -- (1.6,0.55) -- (1.6,0.05) -- (7.6,0.05) -- (7.6,0.55) -- (8.4,0.55);
  \node[hw label, anchor=east] at (0.5,0.3) {$\overline{CS}$/$\overline{OE}$};
  % DATA — 高阻(中间虚线) → 有效数据 → 释放回高阻
  \draw[BookMuted,line width=0.65pt,dashed] (0.6,-0.75) -- (4.6,-0.75);
  \draw[BookMuted,line width=0.65pt,dashed] (7.6,-0.75) -- (8.4,-0.75);
  \draw[BookTeal,line width=0.85pt] (4.6,-0.45) -- (7.6,-0.45) -- (7.6,-1.05) -- (4.6,-1.05) -- cycle;
  \node[hw label, anchor=east] at (0.5,-0.75) {DATA};
  \node[hw label] at (2.6,-0.4) {高阻};
  \node[hw label] at (6.1,-0.75) {数据有效};
  % 区间竖线
  \draw[BookRule,line width=0.45pt,dashed] (1.6,2.15) -- (1.6,-1.3);
  \draw[BookRule,line width=0.45pt,dashed] (4.6,2.15) -- (4.6,-1.3);
  \draw[BookRule,line width=0.45pt,dashed] (7.6,2.15) -- (7.6,-1.3);
  \draw[BookMuted,line width=0.6pt,<->] (1.6,2.0) -- node[hw label, above]{访问时间（等待）} (4.6,2.0);
  \draw[BookMuted,line width=0.6pt,<->] (4.6,2.0) -- node[hw label, above]{CPU 采样窗口} (7.6,2.0);
  \node[hw label, text width=8.6cm] at (4.5,-1.7) {等待状态的目的，是把采样点推迟到数据有效之后，同时保持地址、片选和方向不变。};
\end{pdfdiagram}
\diagramnote{异步 SRAM 读周期：地址先稳定，$\overline{CS}$/$\overline{OE}$ 拉低后存储器需要一段访问时间才把数据驱动到总线上；CPU 只能在数据有效后采样。三条线各有明确的翻转时刻——读时序图读的就是这些边沿，而不是三条平直线。}

\topic{SRAM 写周期的驱动权与写窗口}
写周期与读周期使用同一组地址线、数据线和片选，差别在数据线的驱动方向和写使能窗口。读周期中 $\overline{OE}$ 有效，由 SRAM 驱动数据线，发起者只负责在数据有效后采样；写周期中 $\overline{OE}$ 必须保持无效，SRAM 的输出驱动器关闭，改由发起者把写数据驱动到数据线上——同一组数据线在任何时刻只能有一个驱动者。写使能 $\overline{WE}$ 拉低形成写窗口，规格书围绕它规定三条边界：地址必须在窗口打开前已经稳定，写数据必须在窗口关闭前满足建立时间，窗口关闭后还要继续保持一段保持时间，窗口本身也有最小宽度。这些要求来自同一个事实：存储阵列在 $\overline{WE}$ 回升、窗口关闭的边沿附近把数据线上的值写入选中单元，窗口内任何一次地址或数据翻动，都可能把错误的值写进错误的单元。因此写周期的固定顺序是：先给地址和数据，再打开写窗口，先关窗口，最后才撤数据。

\begin{waveform}[x=0.85cm,y=0.95cm]
  % 区间竖线：x=2 地址稳定, x=4 窗口打开, x=9 窗口关闭(写入边沿), x=9.8 数据释放
  \draw[BookRule, line width=0.45pt, dashed] (2,8.35) -- (2,0.25);
  \draw[BookRule, line width=0.45pt, dashed] (4,8.35) -- (4,0.25);
  \draw[BookRule, line width=0.45pt, dashed] (9,8.35) -- (9,0.25);
  \draw[BookRule, line width=0.45pt, dashed] (9.8,8.35) -- (9.8,0.25);
  % ADDR — 窗口打开前先稳定，窗口关闭后才允许变化
  \draw[BookTeal, line width=0.9pt] (0.5,6.5) rectangle (2,7.2);
  \node at (1.25,6.85) {旧值};
  \draw[BookTeal, line width=0.9pt] (2,6.5) rectangle (12,7.2);
  \node at (7,6.85) {地址稳定};
  \draw[BookMuted, line width=0.65pt, dashed] (12,6.85) -- (13,6.85);
  \node[hw label, anchor=east] at (0.3,6.85) {ADDR};
  % CS_n — 低有效: 地址稳定后拉低
  \draw[BookAccent, line width=0.9pt] (0.5,5.6) -- (2.5,5.6) -- (2.5,4.9) -- (10.5,4.9) -- (10.5,5.6) -- (13,5.6);
  \node[hw label, anchor=east] at (0.3,5.25) {$\overline{CS}$};
  % WE_n — 写窗口
  \draw[BookAccent, line width=0.9pt] (0.5,4.2) -- (4,4.2) -- (4,3.5) -- (9,3.5) -- (9,4.2) -- (13,4.2);
  \node[hw label, anchor=east] at (0.3,3.85) {$\overline{WE}$};
  % OE_n — 全程保持无效
  \draw[BookAccent, line width=0.9pt] (0.5,2.8) -- (13,2.8);
  \node[hw label, anchor=east] at (0.3,2.8) {$\overline{OE}$};
  \node[hw label] at (6.75,2.42) {全程保持无效};
  % DATA — 高阻(中间虚线) → 发起者驱动写数据 → 释放回高阻
  \draw[BookMuted, line width=0.65pt, dashed] (0.5,1.4) -- (4.5,1.4);
  \draw[BookMuted, line width=0.65pt, dashed] (9.8,1.4) -- (13,1.4);
  \draw[BookTeal, line width=0.9pt] (4.5,1.0) -- (9.8,1.0) -- (9.8,1.8) -- (4.5,1.8) -- cycle;
  \node[hw label, anchor=east] at (0.3,1.4) {DATA};
  \node[hw label] at (2.5,1.78) {高阻};
  \node at (7.15,1.4) {发起者驱动：写数据稳定};
  \node[hw label] at (11.4,1.78) {高阻};
  % 顶部标注: 地址先行 / 写窗口
  \draw[BookMuted, line width=0.6pt, <->] (2,7.95) -- node[hw label, above]{地址先稳定} (4,7.95);
  \draw[BookMuted, line width=0.6pt, <->] (4,7.95) -- node[hw label, above]{写窗口} (9,7.95);
  % 底部标注: 建立 / 保持
  \draw[BookMuted, line width=0.6pt, <->] (4.5,0.55) -- node[hw label, below]{数据建立} (9,0.55);
  \draw[BookMuted, line width=0.6pt, <->] (9,0.55) -- node[hw label, below]{数据保持} (9.8,0.55);
\end{waveform}
\diagramnote{异步 SRAM 写周期：地址与片选先稳定，$\overline{OE}$ 全程无效，发起者在写窗口内驱动数据线。存储阵列在 $\overline{WE}$ 回升、窗口关闭的边沿把数据线上的值写入选中单元，因此数据必须在窗口关闭前满足建立时间、关闭后满足保持时间，地址在整个窗口内不能翻动——与读周期对照，差别就在数据线由谁驱动、哪个边沿是真正的写入点。}

\topic{DRAM 需要控制器而不是直接接地址线}
DRAM 的容量来自更密集的存储单元，代价是接口复杂。控制器要先打开某个 bank 的某一行，再读写列；换行前需要预充电；即使没有普通访问，也要周期性刷新；高速 DRAM 还需要训练采样窗口。CPU 看到的是“主存读写”，但实际路径经过控制器、队列、bank 调度、行命中或行冲突。这也是为什么现代 CPU 必须靠 cache 把常用数据放近。

\begin{longtable}{L{0.2\linewidth}L{0.34\linewidth}L{0.36\linewidth}}
\toprule
DRAM 动作 & 硬件含义 & 时延来源 \\
\midrule
Activate & 打开某个 bank 中的一行 & 行译码和读出放大器（感放）建立 \\\tablerowrule
Read/Write & 在已打开行内选择列传输数据 & 列访问、突发传输和总线排队 \\\tablerowrule
Precharge & 关闭当前行，准备访问另一行 & 位线恢复到可再次访问状态 \\\tablerowrule
Refresh & 周期性恢复存储单元电荷 & 控制器必须暂停或调度普通访问 \\
\bottomrule
\end{longtable}

\topic{存储芯片的位宽扩展与容量扩展}
单片存储芯片的位宽和容量都有限，整机用两组方向相反的手段补齐。位宽扩展解决“字不够宽”：把若干片芯片并接同一组地址线、片选和读写控制，每片只接数据线的一个字节车道——四片 x8 SRAM 并接后，一次 32-bit 访问四片同时动作，各自收发 D[7:0]、D[15:8]、D[23:16]、D[31:24]，对外表现为一片 32-bit 宽的 SRAM。容量扩展解决“行不够多”：数据线和低位地址全部并接，新增的高位地址经译码器产生多路互斥片选，任一地址最多选中一片，行数随片数倍增。两种扩展经常同时使用，但边界必须分清：位宽扩展中各片同时被选中、各管一段数据线；容量扩展中各片互斥选中、数据线全部共用。混淆二者的直接后果，是两片同时驱动同一个字节车道的总线争用。

\begin{pdfdiagram}
  \node[hw control, minimum width=2.1cm] (src) at (0.7,0) {同一组地址\\$\overline{CS}$/$\overline{WE}$};
  \node[hw memory, minimum width=2.0cm] (c0) at (4.15,2.7) {x8 SRAM 片 0};
  \node[hw memory, minimum width=2.0cm] (c1) at (4.15,0.9) {x8 SRAM 片 1};
  \node[hw memory, minimum width=2.0cm] (c2) at (4.15,-0.9) {x8 SRAM 片 2};
  \node[hw memory, minimum width=2.0cm] (c3) at (4.15,-2.7) {x8 SRAM 片 3};
  \node[hw state, minimum width=2.2cm, minimum height=6.6cm] (dbus) at (7.9,0) {32-bit 数据\\D[31:0]};
  % 地址/控制共享干线 + 分层支路
  \draw[hw bus, -] (src.east) -- (2.5,0);
  \draw[hw bus, -] (2.5,2.7) -- (2.5,-2.7);
  \draw[hw bus, shorten <=0pt] (2.5,2.7) -- (c0.west);
  \draw[hw bus, shorten <=0pt] (2.5,0.9) -- (c1.west);
  \draw[hw bus, shorten <=0pt] (2.5,-0.9) -- (c2.west);
  \draw[hw bus, shorten <=0pt] (2.5,-2.7) -- (c3.west);
  % 各片各占一个字节车道
  \draw[hw bus] (c0.east) -- node[above,hw label]{D[7:0]} (6.8,2.7);
  \draw[hw bus] (c1.east) -- node[above,hw label]{D[15:8]} (6.8,0.9);
  \draw[hw bus] (c2.east) -- node[above,hw label]{D[23:16]} (6.8,-0.9);
  \draw[hw bus] (c3.east) -- node[above,hw label]{D[31:24]} (6.8,-2.7);
  \node[hw label, text width=2.2cm] at (10.3,0) {容量扩展则靠片选译码增加行数};
\end{pdfdiagram}
\diagramnote{位宽扩展：四片 x8 SRAM 并接同一组地址、$\overline{CS}$ 和 $\overline{WE}$，同时动作、各管一个字节车道，对外是一片 32-bit SRAM。容量扩展方向相反：数据线全部并接，高位地址译码出互斥片选来增加行数，任一地址最多一片被选中。}

\topic{DRAM 单元用电容存电荷，读出后必须重写}
DRAM 的一个 bit 由一个访问晶体管和一个存储电容构成（1T1C）。位线先预充到 $V_{CC}/2$；字线打开晶体管后，电容与位线共享电荷，位线电压略升或略降，读出放大器（感放）判别方向并输出 0/1。两个后果决定了 DRAM 的全部特殊性：其一，读出是破坏性的，感放必须把满摆幅电平经位线写回单元，否则原值丢失；其二，电容持续漏电，电荷在几十毫秒内就衰减到无法判别，控制器必须周期性刷新——逐行打开并重写。规格书里的 Activate、Precharge、Refresh 和刷新周期（典型要求 64 ms 内把全部行刷新一遍），都直接来自这两件事。

\begin{pdfdiagram}
  % 位线
  \draw[hw bus] (0.4,2.2) -- (4.9,2.2);
  \node[hw label, above] at (0.4,2.2) {位线 BL（先预充到 $V_{CC}/2$）};
  % 访问晶体管
  \pic at (1.6,1.55) {hw nmos};
  \draw[hw bus] (1.6,1.97) -- (1.6,2.2);
  \draw[hw bus] (0.4,1.55) -- (0.95,1.55);
  \node[hw label, above] at (0.4,1.55) {字线 WL};
  % 存储电容
  \draw[hw bus] (1.6,1.13) -- (1.6,0.92);
  \draw[hw bus] (1.38,0.92) -- (1.82,0.92);
  \draw[hw bus] (1.38,0.78) -- (1.82,0.78);
  \draw[hw bus] (1.6,0.78) -- (1.6,0.5);
  \pic at (1.6,0.5) {hw gnd};
  \node[hw label, anchor=west] at (2.0,0.85) {存储电容 C};
  % 读出放大器
  \node[hw compute, minimum width=3.0cm] (sa) at (6.5,2.2) {读出放大器（感放）\\判别后满摆幅回写};
  \node[hw label, anchor=west] at (8.15,2.2) {读出的 0/1};
  \node[hw label, text width=9.6cm] at (4.3,-0.75) {打开字线后，电容与位线共享电荷，位线电压轻微上升或下降；感放判别方向输出 0/1，并立刻把满摆幅电平写回单元——读出即重写。};
\end{pdfdiagram}
\diagramnote{1T1C 存储单元与读出放大器。刷新并不是额外功能，而是控制器替所有行周期性重复这个“读出—回写”过程；这也是 DRAM 不能像 SRAM 那样直接把地址线接上去用的原因。}

\topic{DRAM 的行列地址复用与突发传输}
DRAM 的地址引脚不随容量等比增长，靠的是行列地址复用：同一组地址引脚先送行地址、再送列地址。一个 $2^{20}$ 字的阵列直接接地址需要 20 根地址线，复用后 10 根就够——$\overline{RAS}$ 有效时把引脚上的值锁存为行地址，$\overline{CAS}$ 有效时把同一组引脚上的值锁存为列地址。省下的引脚直接降低封装成本，代价是接口从“给地址等数据”变成一串带时序要求的命令。行地址译码后打开整行，整行数据被感放读出并留在行缓冲中；此后的列访问只是从行缓冲里选一段送到数据总线。突发传输利用的正是这一点：给一个起始列地址，DRAM 按内部列计数器连续多拍输出相邻数据，后续各拍不再付出行打开的代价。行缓冲命中是突发便宜的根本原因，也是内存控制器尽量把相邻访问调度到同一行的理由。

\begin{pdfdiagram}
  \node[hw state, minimum width=2.0cm] (addr) at (0.6,0) {地址总线};
  \node[hw compute, minimum width=2.0cm] (mux) at (2.9,0) {行列复用器};
  \node[hw state, minimum width=1.9cm] (rowlatch) at (5.4,0.9) {行地址锁存};
  \node[hw state, minimum width=1.9cm] (collatch) at (5.4,-0.9) {列地址锁存};
  \node[hw compute, minimum width=1.4cm] (rowdec) at (7.7,0.9) {行译码};
  \node[hw compute, minimum width=1.4cm] (coldec) at (7.7,-0.9) {列译码};
  \node[hw memory, minimum width=2.0cm] (bank) at (10.2,0.9) {bank 阵列};
  \node[hw memory, minimum width=2.0cm] (rowbuf) at (10.2,-0.9) {行缓冲};
  \node[hw compute, minimum width=2.0cm] (burst) at (10.2,-2.4) {突发输出\\连续多拍};
  % 地址总线经共享干线分送到行、列锁存
  \draw[hw bus] (addr) -- (mux);
  \draw[hw bus] (mux.east) -- (5.4,0) -- (rowlatch.south);
  \draw[hw bus, shorten <=0pt] (5.4,0) -- (collatch.north);
  \draw[hw return] (5.4,1.55) -- (rowlatch.north);
  \node[hw label] at (5.4,1.85) {$\overline{RAS}$};
  \draw[hw return] (5.4,-1.55) -- (collatch.south);
  \node[hw label] at (5.4,-1.85) {$\overline{CAS}$};
  % 行路径与列路径
  \draw[hw bus] (rowlatch) -- node[above,hw label]{行} (rowdec);
  \draw[hw bus] (collatch) -- node[above,hw label]{列} (coldec);
  \draw[hw bus] (rowdec) -- (bank);
  \draw[hw bus] (coldec) -- (rowbuf);
  \draw[hw bus] (bank.south) -- node[right,hw label]{整行读入} (rowbuf.north);
  \draw[hw bus] (rowbuf.south) -- node[right,hw label]{按列连读} (burst.north);
  \node[hw label, text width=10.2cm] at (5.35,-3.5) {同一组地址引脚在 $\overline{RAS}$、$\overline{CAS}$ 两拍分别送出行、列地址；整行先读入行缓冲，列访问只是在缓冲里选段，行命中后的连续突发因此不必重开行。};
\end{pdfdiagram}
\diagramnote{行列地址复用与突发传输。$\overline{RAS}$ 拍锁存行地址并打开整行，$\overline{CAS}$ 拍锁存列地址；行缓冲命中后，连续列地址直接从行缓冲多拍输出，比重开一行便宜得多。}

行打开、列传输和预充电之间的最小间隔由一组时序参数规定，规格书把它们写成以时钟拍为单位的数字。DDR3-1600 的内部时钟为 800 MHz，一拍 1.25 ns；1600 MT/s 的数据速率来自同一时钟的双沿传输，时序参数仍按时钟拍计。常见标注 11-11-11-28 就是下表四个数字。

\begin{longtable}{L{0.12\linewidth}L{0.44\linewidth}L{0.16\linewidth}L{0.18\linewidth}}
\toprule
参数 & 含义 & DDR3-1600 & 纳秒换算 \\
\midrule
$t_{RCD}$ & Activate 打开行之后，到第一个列读/写命令的最小间隔 & 11 拍 & 13.75 ns \\\tablerowrule
$t_{CL}$ & 列读命令发出后，到第一个数据出现在总线上（CAS 延迟） & 11 拍 & 13.75 ns \\\tablerowrule
$t_{RP}$ & Precharge 关闭当前行之后，到允许再次 Activate 的最小间隔 & 11 拍 & 13.75 ns \\\tablerowrule
$t_{RAS}$ & 一次 Activate 到允许 Precharge 的最小行打开时间 & 28 拍 & 35 ns \\
\bottomrule
\end{longtable}

换算很直接：拍数乘 1.25 ns。同一行内连续读，列延迟 13.75 ns 后就能紧跟突发数据；访问另一行则要先等 $t_{RP}$ 再重新 Activate，一次行冲突至少多花 $t_{RP}+t_{RCD}=22$ 拍，约 27.5 ns。这个数字就是内存调度器在乎行命中、软件在乎访问局部性的原因。

\topic{SDRAM 的突发长度与多 bank 交错}
Activate 打开一行之后，控制器可以对同一行连发多个列读写命令；每个列命令传的也不是一个数据，而是按约定的突发长度（burst length）连续传一串——DDR3 固定为 8（BL8），也允许切成 4。DDR 的“双倍数据率”指时钟的上升沿和下降沿各传一次数据，而不是把时钟翻倍：DDR3-1600 的内部时钟是 800 MHz，数据速率 1600 MT/s，所以一次 BL8 突发占 4 个时钟拍、交 8 个数据，这期间数据总线被这一串独占。

bank 是 DRAM 内部可独立 Activate、独立 Precharge 的子阵列，各有自己的行缓冲。一个 bank 在等 $t_{RCD}$、等数据的时候，命令总线可以对另一个 bank 发命令，把后者的行打开时间藏进前者的传输里，这就是多 bank 交错。沿用上一张表的 DDR3-1600 参数（1 拍 1.25 ns，$t_{RCD}=t_{CL}=t_{RP}=11$ 拍，$t_{RAS}=28$ 拍），比较“读两个不同行”在单 bank 与双 bank 下的时间轴：

\begin{lstlisting}
// 注释（推进）：按行追踪数据来源、经过的部件和最终写入目标。
// 注释（边界）：只有标出的时钟沿、阶段或异常入口会改变体系结构可见状态。
单 bank，两次访问落在同一 bank 的不同行:
  拍  0  ACT bank0 row0
  拍 11  RD  bank0 col0    (t_RCD 到期)
  拍 22  数据出现在总线，BL8 占用拍 22-25
  拍 28  PRE bank0         (t_RAS 到期才允许关行)
  拍 39  ACT bank0 row1    (t_RP 到期)
  拍 50  RD  bank0 col1    (t_RCD 到期)
  拍 61  第二组数据出现，共等 61 拍

双 bank 交错，两行分别在 bank0、bank1:
  拍  0  ACT bank0
  拍  1  ACT bank1         (不同 bank，行打开并行进行)
  拍 11  RD  bank0         (数据占拍 22-25)
  拍 15  RD  bank1         (数据占拍 26-29，恰好接上)
  拍 26  第二组数据出现，共等 26 拍
\end{lstlisting}

第二组数据从拍 61 提前到拍 26，省 35 拍、约 43.75 ns。省下的不是任何一条时序参数——单 bank 必须串行等完 $t_{RP}+t_{RCD}$ 才能再读，双 bank 让 bank1 的 Activate 与 bank0 的全过程并行，这段等待被完全覆盖；拍 15 才发第二条读命令，是因为数据总线仍由两个 bank 共享，bank0 的突发占着拍 22--25，bank1 的数据最早只能接在拍 26。这就是内存条标注 bank 数、控制器把相邻地址散到不同 bank、调度器重排请求的原因：行命中决定一次访问多快，bank 并行决定等待能不能被藏起来。这个例子是教学模型，为突出 bank 并行的收益，忽略了 $t_{RRD}$、$t_{FAW}$ 等对连续 Activate 的命令间约束。

\begin{longtable}{L{0.22\linewidth}L{0.32\linewidth}L{0.32\linewidth}}
\toprule
对照点 & 单 bank 顺序访问 & 双 bank 交错 \\
\midrule
第二组数据出现 & 拍 61 & 拍 26（省 35 拍，约 43.75 ns） \\\tablerowrule
数据总线占用 & 拍 22-25、61-64，中间大量气泡 & 拍 22-29 连续无气泡 \\\tablerowrule
关键等待 & 换行时 $t_{RP}+t_{RCD}$ 串行暴露 & 被另一 bank 的并行 Activate 覆盖 \\\tablerowrule
检查点 & PRE 与 ACT 都卡在同一个 bank 上 & 命令错开，数据在共享总线上排队衔接 \\
\bottomrule
\end{longtable}

\topic{组相联 cache 的替换策略是精度与代价的取舍}
新行取回时组内已无空路，牺牲哪一路？直接映射没有这个问题，每个地址只有唯一槽位；组相联把它交给了替换策略。

精确 LRU（最久未使用）记录组内各路的年龄顺序，牺牲最老的一路，命中率最好，代价随路数急剧膨胀。2 路组相联时每组只需 1 位：这一位指出哪一路更久未用，每次命中取反即可。4 路要区分完整的年龄顺序，共 $4!=24$ 种排列，每组需 $\lceil\log_2 24\rceil=5$ 位，且每次命中都要重排年龄；8 路则需 $\lceil\log_2 40320\rceil=16$ 位。伪 LRU 把 4 路组织成两层二叉树，3 个内部节点各 1 位，每组 3 位，命中时沿路径翻转节点位；它选出的不保证是最老路，但命中率损失通常很小。FIFO 只记进入顺序，命中时不更新任何状态；随机替换连状态位都不要。

高组相联之所以普遍放弃精确 LRU，原因有两头：状态位按 $n!$ 增长还是小事，更难办的是每次命中都要更新年龄，路数越多更新逻辑的组合路径越长，直接拖慢命中时间；而 8 路、16 路的组里，“精确选最老”换来的命中率收益早已递减，随机替换与 LRU 的差距常常小到测量噪声量级。工程上的结论因此很清楚：2 路用 1 位精确 LRU，4 路用 3 位伪 LRU，更高组相联用伪 LRU 或随机。

\begin{longtable}{L{0.15\linewidth}L{0.25\linewidth}L{0.20\linewidth}L{0.28\linewidth}}
\toprule
策略 & 每组状态位（4 路） & 命中时是否更新 & 关键问题 \\
\midrule
精确 LRU & 5 位（$4!=24$ 种年龄序） & 每次命中重排年龄 & 路数增多后状态位和更新路径急剧膨胀 \\\tablerowrule
伪 LRU & 3 位（两层树节点） & 沿树路径翻转节点位 & 选出的不保证是最老路，命中率略降 \\\tablerowrule
FIFO & 2 位（轮流指针） & 不更新 & 不看访问热度，热行可能被轮到 \\\tablerowrule
随机 & 0 位 & 不更新 & 实现最简，大组相联下命中率损失很小 \\
\bottomrule
\end{longtable}

\topic{硬件预取与软件预取都是拿带宽换命中}
替换策略决定牺牲谁，预取则试图让牺牲更少发生：顺序扫描数组时，用完一条 cache 行，下一行几乎马上要用，空间局部性让“将来要访问什么”变得可预测。硬件预取器监视 miss 流的地址规律：next-line 预取在用到第 $i$ 行时把第 $i+1$ 行取回；步长预取识别出固定步长后沿步长提前取回；预取流通常在 4 KiB 页边界停下——跨页地址可能未映射，为一次猜测触发缺页代价太大。

软件预取由指令显式给出提示。x86-64 提供 \code{prefetcht0}、\code{prefetchnta} 等预取指令：它们不改变任何架构可见状态，地址无效时也不触发缺页异常，只是提示硬件“这个地址很快会被读”。编译器或手写汇编常在循环里提前若干次迭代发出提示：

\begin{lstlisting}
; rdi 指向当前元素, rsi 是数组结尾, 每元素 8 字节
scan_loop:
    prefetcht0 [rdi + 512]   ; 提前 8 行(64 B 行)提示
    add     rax, [rdi]
    add     rdi, 8
    cmp     rdi, rsi
    jb      scan_loop
\end{lstlisting}

预取不是免费的：取错行会占住 cache 槽位、挤掉有用数据，每次预取本身也消耗内存带宽，过激的预取反而把真正需要的访问排在后面。可以粗算一笔账：顺序扫描大数组，行 64 B、元素 8 B，一行装 8 个元素；miss 时延按 200 拍、每个元素处理 4 拍计。无预取时每行一次 miss，每行花 $200+32=232$ 拍；预取把“取下一行”与“处理当前行”重叠，若提前量达到 $200\div 32\approx 7$ 行，处理完当前行时下一行已经在 cache 里，稳态每行只剩 32 拍，miss 被完全隐藏，吞吐提高约 $232/32\approx 7.3$ 倍。若预取只能提前一行，32 拍的处理盖不住 200 拍的取回，每行仍暴露约 168 拍。提前量必须匹配“处理一行的耗时”与“取回一行的耗时”之比，这正是预取器设计的核心数字。

\begin{longtable}{L{0.18\linewidth}L{0.34\linewidth}L{0.34\linewidth}}
\toprule
预取方式 & 触发时机 & 关键问题 \\
\midrule
next-line 硬件预取 & 用到第 $i$ 行时取第 $i+1$ 行 & 顺序流效果好，随机访问纯属浪费 \\\tablerowrule
步长硬件预取 & 识别出固定地址步长后沿步长取 & 步长误判会持续污染 cache \\\tablerowrule
软件预取指令 & 程序在用到数据之前显式提示 & 提示距离要折算成拍数，地址须确实会用 \\\tablerowrule
跨页停止 & 预取流在 4 KiB 页边界停下 & 猜测性访问不得触发缺页副作用 \\
\bottomrule
\end{longtable}

\topic{上电自检的走步与棋盘格各查一类故障}
主存焊上板不等于可信：某位固定读成 0 或 1（stuck-at）、相邻数据线或地址线短路、写入后状态维持不住（数据保持故障），都可能让“写一个值再读回来”这种最简单的测试侥幸通过——它只验证了大部分位在此刻、对这组值是可写的。上电自检用图案化的写入-读回序列把故障类别分开：走步（walking 1/0）每次只让一位与众不同，棋盘格（checkerboard）让相邻位恒取反值。

走步 1 先写全 0 背景，再对同一单元依次写 \code{0x01}、\code{0x02}、……、\code{0x80}，每步读回比对，并复查其余单元仍是背景值。某位固定 0 会在轮到它走 1 时露馅，固定 1 会在背景检查中露馅；两位短路时走步值会变形——写 \code{0x01} 可能读回 \code{0x03}，直接指认相邻两位被短接。走步 0 把背景换成全 1、逐位走 0，对称覆盖另一半故障方向。棋盘格把 \code{0x55}（0101 0101）写满读回，再换成 \code{0xAA}（1010 1010）：任何相邻两位恒取反，相邻短路或位间耦合会让图案当场变形；保持故障则靠“写入后停一段再读”暴露。这些图案便宜、确定、能在无操作系统的裸机上运行，查不出所有故障，但恰好覆盖焊接与走线最常出的几类，因此是板级 bring-up 的固定节目。

以 4 字节小内存（addr0--addr3，8-bit 字节）为例，一次走步 1 加棋盘格的完整步骤如下。

\begin{longtable}{L{0.08\linewidth}L{0.38\linewidth}L{0.18\linewidth}L{0.24\linewidth}}
\toprule
步骤 & 操作 & 预期结果 & 检查点 \\
\midrule
1 & addr0--addr3 全写 \code{0x00}，逐字节读回 & 全部 \code{0x00} & 无固定 1，背景可写可读 \\\tablerowrule
2 & addr0 依次写 \code{0x01}、\code{0x02}、……、\code{0x80}，每步读回 & 读回恒等于写入 & 每个数据位都能独立置 1 \\\tablerowrule
3 & 步骤 2 每步之后读 addr1--addr3 & 仍为 \code{0x00} & 无地址串扰、无相邻线短路 \\\tablerowrule
4 & addr1--addr3 重复步骤 2、3 & 同上 & 逐字节覆盖全部单元 \\\tablerowrule
5 & 全写 \code{0x55} 读回，再全写 \code{0xAA} 读回 & 图案不变形 & 相邻位恒反，查相邻短路 \\\tablerowrule
6 & 背景改 \code{0xFF}，逐位走 0（\code{0xFE}、\code{0xFD}、……） & 读回恒等于写入 & 对称覆盖固定 0 故障 \\
\bottomrule
\end{longtable}

\topic{校验位与 ECC 把位错从灾难降级为事件}
主存的位会错：宇宙射线、电压毛刺、保持时间失守都可能翻转一个 bit。最便宜的对策是奇偶校验——每字节加 1 个校验位，使连同校验位在内的 1 的个数恒为偶数（偶校验）。数据 \code{1011 0001} 已有 4 个 1，校验位记 0；读回时若某位翻转，1 的个数变奇，硬件立刻知道出错了。奇偶校验的边界同样明确：只能发现奇数个位错，两位同时翻转就漏检；即使发现，也不知道错在哪一位，只能报错，不能纠正。

汉明码用多个校验位换取定位能力。(7,4) 汉明码把 4 个数据位和 3 个校验位排成 7 位，位置从 1 编号；校验位放在编号为 2 的幂的位置（1、2、4），数据位占据 3、5、6、7。把位置编号写成 3 位二进制，第 $i$ 个校验位负责所有编号第 $i$ 位为 1 的位置。每个校验组各做一次偶校验，哪几组失败，其编号拼起来恰好就是出错位置的编号——这个把“哪些组错”翻译成“哪里错”的数称为校正子（syndrome）。

\begin{longtable}{L{0.12\linewidth}L{0.16\linewidth}L{0.28\linewidth}L{0.28\linewidth}}
\toprule
位置编号 & 存放内容 & 参与校验的校验位 & 说明 \\
\midrule
1 & $p_1$ & $p_1$ 自身 & 编号 001，最低位为 1 \\\tablerowrule
2 & $p_2$ & $p_2$ 自身 & 编号 010 \\\tablerowrule
3 & $d_1$ & $p_1$、$p_2$ & 编号 011 \\\tablerowrule
4 & $p_4$ & $p_4$ 自身 & 编号 100 \\\tablerowrule
5 & $d_2$ & $p_1$、$p_4$ & 编号 101 \\\tablerowrule
6 & $d_3$ & $p_2$、$p_4$ & 编号 110 \\\tablerowrule
7 & $d_4$ & $p_1$、$p_2$、$p_4$ & 编号 111 \\
\bottomrule
\end{longtable}

算一遍：数据位（位 3、5、6、7）依次为 1、0、1、1。$p_1$ 覆盖位 1、3、5、7，其中数据 1、0、1 共两个 1，$p_1=0$；$p_2$ 覆盖位 2、3、6、7，数据 1、1、1 共三个 1，$p_2=1$；$p_4$ 覆盖位 4、5、6、7，数据 0、1、1 共两个 1，$p_4=0$。发出码字（位 1 至 7）为 0 1 1 0 0 1 1。若位 6 在途中翻成 0：$p_1$ 组（0,1,0,1）共两个 1，通过；$p_2$ 组（1,1,0,1）共三个 1，失败；$p_4$ 组（0,0,0,1）共一个 1，失败。校正子按 $p_4p_2p_1$ 读出 $(1,1,0)=6$，直接指向位 6，翻转回来即完成纠正。

实用的 ECC 内存把同一思想做到 64 位数据加 8 位校验、72 位宽，实现单错纠正、双错检测（SEC-DED）。贵在三处：阵列宽了 1/8，每条内存要多挂芯片；控制器对每次读做译码纠正，对部分写先读-改-写再重算校验，多一层逻辑和时延；平台还要为纠错路径做完整验证。对长时间运行、数据不能错的机器这是便宜的保险，桌面机则普遍省掉——这也是普通内存条与 ECC 内存条不能随意混插的原因。

\section{三、存储层次、地址转换和设备驱动 trace（执行轨迹）}

一次普通 LOAD 可能先经过 TLB 查找，再查 L1 cache，未命中后查 L2/L3，最后进入内存控制器；一次设备接收数据可能由驱动准备描述符，设备 DMA 写缓冲区，cache 一致性或维护操作让 CPU 看见新数据，中断再把完成事件带回来。

\begin{longtable}{L{0.18\linewidth}L{0.4\linewidth}L{0.32\linewidth}}
\toprule
访问层次 & 常见命中路径 & 失败或慢路径 \\
\midrule
寄存器 & 直接从寄存器堆读写 & 端口冲突、旁路失败或流水线停顿 \\\tablerowrule
L1 cache & 近处 SRAM 命中，几个周期内返回 & L1 miss，向 L2 请求整条 cache line \\\tablerowrule
L2/L3 cache & 片上较大 cache 命中 & 最后级 miss，访问内存控制器和 DRAM \\\tablerowrule
TLB & 地址转换缓存命中 & page walk、页错误或权限异常 \\\tablerowrule
DRAM & 控制器调度行命中和突发传输 & 行冲突、刷新、排队、ECC 错误 \\\tablerowrule
设备/MMIO & 设备寄存器响应读写 & 背压、超时、错误状态或中断完成 \\
\bottomrule
\end{longtable}

\topic{TLB 把地址转换也纳入近处缓存}
虚拟地址或线性地址到物理地址的转换若每次都查页表，访存会变得很慢。TLB 可以理解为“地址转换的 cache”：它保存近期虚拟页到物理页框的映射及权限位。LOAD/STORE 形成虚拟地址后，先用虚拟页号查 TLB；命中时得到物理页框和权限信息，再与页内偏移组合成物理地址；未命中时，硬件 page walker 或软件异常处理会去内存里的页表结构查找。

\begin{lstlisting}
virtual address = [ virtual_page_number | page_offset ]
TLB hit:
  physical_address = TLB[virtual_page_number].frame | page_offset
TLB miss:
  walk page tables or enter miss handler
\end{lstlisting}

\begin{longtable}{L{0.18\linewidth}L{0.36\linewidth}L{0.36\linewidth}}
\toprule
TLB 字段 & 作用 & 关键问题 \\
\midrule
虚拟页号 tag & 判断当前地址是否命中某个转换项 & 进程切换、地址空间标识和刷新规则是否清楚 \\\tablerowrule
物理页框 & 与页内偏移组合成物理地址 & 页大小、对齐和物理地址宽度是否一致 \\\tablerowrule
权限位 & 读写、执行、用户/特权等访问属性 & 访问失败是权限异常，不是普通 cache miss \\\tablerowrule
valid/global & 表示转换是否有效、是否跨进程保留 & 页表修改后如何让旧 TLB 项失效 \\
\bottomrule
\end{longtable}

TLB miss 和 cache miss 都是“近处没有”，但语义不同。cache miss 只是数据不在近处，通常可以向下级取回；TLB miss 可能只是转换缓存缺失，也可能说明页不存在或权限不允许。若页表项不存在，硬件不能继续访问内存数据；若权限不允许，访问必须转入异常。这个差别解释了为什么地址转换、保护和异常在系统硬件中是同一条路径的不同分支。

\topic{多核和 DMA 让 cache 一致性成为硬件约定}
单核教学机里，cache 可以先理解为 CPU 与主存之间的近处副本。多核或 DMA 加入后，问题变成“多个观察者如何看见同一地址的值”。一个核心的 cache 中可能有某条 line 的旧副本，另一个核心或设备刚刚写入了新值；如果没有一致性协议、cache 维护或不可缓存映射，CPU 读到的就可能不是系统真实最新值。

\begin{longtable}{L{0.2\linewidth}L{0.36\linewidth}L{0.34\linewidth}}
\toprule
场景 & 风险 & 常见硬件/软件边界 \\
\midrule
核心 A 写，核心 B 读 & B 的 cache 中可能有旧行 & cache coherence、失效消息、共享/独占状态 \\\tablerowrule
CPU 写描述符，设备读 & 描述符可能还在 CPU cache 或写缓冲中 & coherent DMA、flush、写屏障 \\\tablerowrule
设备写缓冲区，CPU 读 & CPU cache 中可能有旧数据 & invalidate、coherent 互连、不可缓存映射 \\\tablerowrule
MMIO 状态寄存器 & cache 命中会隐藏设备新状态 & 设备内存属性、禁止普通缓存 \\
\bottomrule
\end{longtable}

一致性协议细节可以很深，但教学阶段先抓住三件事。第一，同一地址可能在多个近处缓存中有副本。第二，写入必须让其他观察者的旧副本失效或更新。第三，设备并不总是自动参与 CPU cache 一致性；若平台不支持 coherent DMA，驱动必须在交出缓冲区前后做 cache 维护。否则 DMA 完成中断到达后，CPU 仍可能读到旧 cache line。

\topic{MESI 用四个状态约定谁拿着最新值}
这套状态存在的意义，是让任何时刻、任何地址在所有观察者眼里只有一份“最新值”。想写某一行，核必须先成为唯一持有者：从 S 升级要广播失效，把其他核的副本打成 I；只有处于 M 或 E 时才能直接改。想读可以共享：多个核同时持有 S，读到的都是同一个干净副本，但共享期间谁也不能直接写。于是两个核不会各自拿着不同的新值——要么一个核独占、其余无效，要么大家一起只读。参与一致性互连的 DMA 设备也被同一套状态约束；不参与的平台，驱动就要用 cache clean/invalidate 手工扮演“写回”和“失效”的角色。

\begin{pdfdiagram}
  \node[hw state, minimum width=1.9cm] (stI) at (1.6,3.2) {I 无效};
  \node[hw state, minimum width=1.9cm] (stS) at (9.6,3.2) {S 共享};
  \node[hw state, minimum width=1.9cm] (stE) at (1.6,0.5) {E 独占};
  \node[hw state, minimum width=1.9cm] (stM) at (9.6,0.5) {M 已修改};
  % 本核读：I 分裂为 S 或 E
  \draw[hw return, ->] (stI.east) -- node[above, hw label] {本核读：发现他人副本} (stS.west);
  \draw[hw return, ->] (stI.south) -- (stE.north);
  \node[hw label, anchor=west] at (1.7,1.85) {本核读：无其他副本};
  % 本核写与侦听他人读：S 与 M 之间
  \draw[hw return, <->] (stS.south) -- (stM.north);
  \node[hw label, anchor=east] at (9.5,2.2) {本核写：广播失效 $\rightarrow$ M};
  \node[hw label, anchor=east] at (9.5,1.5) {他人读：写回 $\rightarrow$ S};
  % 本核写：E 已独占，直接改
  \draw[hw return, ->] (stE.east) -- node[above, hw label] {本核写：已独占，直接改} (stM.west);
  % 侦听他人写：副本失效
  \draw[hw return, ->] (stS.north) -- (9.6,3.95) -- (1.6,3.95) -- (stI.north);
  \node[hw label, above] at (5.6,3.95) {侦听他人写：副本失效};
\end{pdfdiagram}
\diagramnote{MESI 一致性状态图（教学简化版）。本核读把 I 变成 E（无其他副本）或 S（他人有副本）；本核写必须先独占：S 到 M 要广播失效，E 到 M 直接改；总线侦听到他人读时 M 写回并退到 S；侦听到他人写时副本失效（E、M 同样失效，M 先写回）。}

\topic{案例：MESI 在两个核之间的完整走查}
状态图只有放进具体访问序列才有意义。下面用两个核、一个地址 X 把上一专题的状态机完整走一遍。设初始时核 A、核 B 的 cache 中 X 行均为 I（无效），内存持有 X 的最新值。先说明一个细节：按严格 MESI，核 A 第一次读入 X 时若没有发现任何其他副本，应记为 E（独占）；本走查假设 X 此前已被第三个核读过，或采用“读入即共享”的教学约定，因此第一步记为 S。这不影响后续步骤的转换逻辑。

\begin{longtable}{L{0.06\linewidth}L{0.24\linewidth}L{0.09\linewidth}L{0.09\linewidth}L{0.42\linewidth}}
\toprule
步骤 & 动作 & 核 A & 核 B & 互连上发生的事 \\
\midrule
0 & 初始 & I & I & X 的最新值在内存中 \\\tablerowrule
1 & 核 A 读 X & S & I & A 的读 miss，经互连从内存取回整行，记为共享副本 \\\tablerowrule
2 & 核 B 读 X & S & S & B 的读 miss，从内存取行；A 侦听到这次读，保持 S，两核持有同一份干净副本 \\\tablerowrule
3 & 核 A 写 X & M & I & A 先广播失效消息，B 的副本作废（S $\rightarrow$ I）；确认后 A 独占修改本行（S $\rightarrow$ M），内存中的 X 从此过期 \\\tablerowrule
4 & 核 B 读 X & S & S & B 的读 miss；A 侦听到读，把脏行写回内存（或直接转发给 B），自己降为 S；B 得到最新数据的共享副本 \\
\bottomrule
\end{longtable}

逐步看两处关键转换。第 3 步是“写之前先独占”：A 不能只改自己手里的 S 副本，否则 B 会继续读到旧值，同一地址就出现了两个“最新值”；广播失效把其他副本全部打成 I 之后，A 才升级成 M 并修改。第 4 步是“最新值跟着 M 走”：此时内存里的 X 是旧数据，B 的读不能由内存直接回答，必须由持有 M 行的 A 先把数据交出来——写回内存或 cache 之间直接转发——之后两核各持 S，内存恢复最新。

这张表就是一致性协议的 trace：每一行对应状态图上的一条边，每一步结束都满足同一条不变式——X 的最新值在全局唯一，要么在某核的 M 行里（其余全为 I，内存过期），要么在所有 S 副本与内存里保持一致。调试多核程序时若怀疑数据不一致，最有效的办法正是把可疑的访问序列按这种方式列成表，逐步检查哪一步破坏了唯一性。

\topic{store buffer 和屏障把“执行完成”与“可见”分开}
STORE 对流水线来说可以很早“执行完成”：地址和数据已经形成，写事务进入 store buffer，后续独立指令可以继续执行。但这个 STORE 什么时候对其他核心、DMA 或设备可见，要看 store buffer 排队、cache 状态、内存类型和屏障规则。多核内存模型本身超出本书范围，这里只需要与 DMA 直接相关的一条规则：普通 RAM 写可以合并和延迟；MMIO 写通常要求更强顺序；DMA 描述符写在通知设备前必须已经对设备可见。

\begin{lstlisting}
// 注释（推进）：逐行追踪发起者、所有权和可见性；请求被接受不等于事务完成。
// 注释（边界）：以采样沿、状态位或 completion 为准，再允许消费者使用结果。
prepare DMA descriptor:
  descriptor.addr = buffer
  descriptor.len  = length
  memory_barrier()
  MMIO[DEVICE_DOORBELL] = descriptor_index
\end{lstlisting}

这里的屏障不是“让 CPU 慢一点”的软件仪式，而是要求硬件在门铃 MMIO 写对设备可见之前，先让描述符内容按平台规则可见。若顺序反了，设备可能看到门铃，却读到旧描述符或半写入描述符。许多难复现的设备错误，本质上就是这个可见性边界没有写清楚。

\begin{longtable}{L{0.18\linewidth}L{0.38\linewidth}L{0.34\linewidth}}
\toprule
对象 & 可见性要求 & 失败表现 \\
\midrule
普通数据 & 由语言、编译器和内存模型共同约束 & 多核读到旧值或乱序观察 \\\tablerowrule
锁变量 & 获取/释放顺序必须约束临界区数据 & 进入临界区后仍看见旧数据 \\\tablerowrule
DMA 描述符 & 通知设备前必须对设备可见 & 设备搬错地址或长度 \\\tablerowrule
MMIO 门铃 & 前面的配置写必须先到设备接口 & 设备按旧配置启动 \\\tablerowrule
中断完成 & completion 到达后数据必须已可读取或可同步 & 中断处理读到旧缓冲区 \\
\bottomrule
\end{longtable}

\topic{案例：一个内存屏障挡住什么}
屏障的作用可以用一个经典的双核序列看清。两个核共享变量 x 和 y，初值都为 0：核 A 先把 1 写入 x，再读 y 到 r1；核 B 先把 1 写入 y，再读 x 到 r2。直觉上，两次写和两次读无论怎样交错，总有一个读应该看到对方的新值 1。但在带 store buffer 的真实机器上，r1=0 且 r2=0 的结果确实可能出现。

\begin{lstlisting}
core A                    core B
mov qword ptr [x], 1      mov qword ptr [y], 1
mov rax, [y]              mov rbx, [x]
\end{lstlisting}

原因就在上一专题的可见性划分：\code{x=1} 对核 A 来说“执行完成”只表示写事务进入了 A 的 store buffer，并不意味着其他核已经能看见它。两个核的 store 都可能还排在各自的 store buffer 里，两个 load 就已经从内存（或对方尚未更新的 cache 副本）里取回了 0。

\begin{longtable}{L{0.08\linewidth}L{0.30\linewidth}L{0.30\linewidth}L{0.22\linewidth}}
\toprule
次序 & 核 A & 核 B & 内存中可见的 x / y \\
\midrule
1 & x=1 进入 A 的 store buffer & --- & 0 / 0 \\\tablerowrule
2 & --- & y=1 进入 B 的 store buffer & 0 / 0 \\\tablerowrule
3 & r1=y，读到 0 & --- & 0 / 0 \\\tablerowrule
4 & --- & r2=x，读到 0 & 0 / 0 \\\tablerowrule
5 & store buffer 排空，x=1 全局可见 & store buffer 排空，y=1 全局可见 & 1 / 1 \\
\bottomrule
\end{longtable}

\begin{pdfdiagram}
  % 表头
  \node[hw label] at (1.6,0.62) {store buffer A};
  \node[hw label] at (4.0,0.62) {内存};
  \node[hw label] at (6.4,0.62) {store buffer B};
  % ① 初始
  \node[hw label, anchor=east] at (-0.65,0) {① 初始};
  \node[hw state, minimum width=0.95cm] (a1) at (0,0) {核 A};
  \node[hw control, minimum width=1.0cm] (ba1) at (1.6,0) {（空）};
  \node[hw memory, minimum width=2.0cm] (m1) at (4.0,0) {x = 0，y = 0};
  \node[hw control, minimum width=1.0cm] (bb1) at (6.4,0) {（空）};
  \node[hw state, minimum width=0.95cm] (b1) at (8.0,0) {核 B};
  % ② 写滞留
  \node[hw label, anchor=east] at (-0.65,-1.5) {② 写滞留};
  \node[hw state, minimum width=0.95cm] (a2) at (0,-1.5) {核 A};
  \node[hw control, minimum width=1.0cm] (ba2) at (1.6,-1.5) {x = 1};
  \node[hw memory, minimum width=2.0cm] (m2) at (4.0,-1.5) {x = 0，y = 0};
  \node[hw control, minimum width=1.0cm] (bb2) at (6.4,-1.5) {y = 1};
  \node[hw state, minimum width=0.95cm] (b2) at (8.0,-1.5) {核 B};
  \draw[hw bus] (a2) -- node[above=1pt,hw label,fill=white,inner sep=1pt]{x=1} (ba2);
  \draw[hw bus] (b2) -- node[above=1pt,hw label,fill=white,inner sep=1pt]{y=1} (bb2);
  % ③ 读旧值
  \node[hw label, anchor=east] at (-0.65,-3.0) {③ 读旧值};
  \node[hw state, minimum width=0.95cm] (a3) at (0,-3.0) {核 A};
  \node[hw control, minimum width=1.0cm] (ba3) at (1.6,-3.0) {x = 1};
  \node[hw memory, minimum width=2.0cm] (m3) at (4.0,-3.0) {x = 0，y = 0};
  \node[hw control, minimum width=1.0cm] (bb3) at (6.4,-3.0) {y = 1};
  \node[hw state, minimum width=0.95cm] (b3) at (8.0,-3.0) {核 B};
  \draw[hw bus] (3.6,-3.35) -- (3.6,-3.8) -- node[pos=0.5,above=1pt,hw label]{r1 = y = 0} (0,-3.8) -- (a3.south);
  \draw[hw bus] (4.4,-3.35) -- (4.4,-4.05) -- node[pos=0.5,below=1pt,hw label]{r2 = x = 0} (8.0,-4.05) -- (b3.south);
  % ④ 屏障后排空
  \node[hw label, anchor=east] at (-0.65,-4.9) {④ 屏障后排空};
  \node[hw state, minimum width=0.95cm] (a4) at (0,-4.9) {核 A};
  \node[hw control, minimum width=1.0cm] (ba4) at (1.6,-4.9) {（空）};
  \node[hw memory, minimum width=2.0cm] (m4) at (4.0,-4.9) {x = 1，y = 1};
  \node[hw control, minimum width=1.0cm] (bb4) at (6.4,-4.9) {（空）};
  \node[hw state, minimum width=0.95cm] (b4) at (8.0,-4.9) {核 B};
  \draw[hw bus] (ba4) -- node[above=1pt,hw label,fill=white,inner sep=1pt]{x=1} (m4);
  \draw[hw bus] (bb4) -- node[above=1pt,hw label,fill=white,inner sep=1pt]{y=1} (m4);
  \node[hw label] at (4.0,-5.75) {mfence 之后：写先全局可见，后续 load 才允许执行};
\end{pdfdiagram}
\diagramnote{store buffer 双核序列的四张快照。x=1、y=1 先滞留在各自核的 store buffer（②），两个 load 直接从内存读回旧值 0（③），于是 r1=0 且 r2=0 成立；mfence 强制 store buffer 排空、写全局可见之后才放行后续 load（④），这个结果即被禁止。四张快照与上表第 1、2、5 步前后的状态一一对应。}

这个结果不违反任何一条单核规则——每个核自己看来程序都按序执行了——但它违反了“所有核看到同一个全局顺序”的直觉。x86 的 \code{mfence} 指令正是为此存在：它强制本核的 store buffer 排空之后，后续的 load 才允许执行，也就是说本核此前的写必须先对其他核可见。

\begin{lstlisting}
core A                    core B
mov qword ptr [x], 1      mov qword ptr [y], 1
mfence                    mfence
mov rax, [y]              mov rbx, [x]
\end{lstlisting}

加入屏障后，r1=0 且 r2=0 被禁止。推理只需直觉：若 A 读到了 y=0，说明 B 的 y=1 在那一刻尚未对 A 可见；而 B 的 \code{mfence} 保证 y=1 全局可见之后 B 才会执行 r2=x，到那时 A 的 x=1 早已因 A 自己的 \code{mfence} 而对所有核可见，所以 B 读 x 必得 1。屏障挡住的正是“写还留在 store buffer 里、读却已经出发”这段窗口；它没有让写变得更快，只是把可见性顺序限定。

\topic{描述符环把 DMA 变成持续事务流}
高速设备通常不会每次只处理一个缓冲区。网卡、磁盘控制器、USB 控制器常使用描述符环或队列：CPU 准备一批描述符，设备按头尾指针消费，完成后写回状态或 completion 队列，并用中断或轮询通知 CPU。这样可以减少 MMIO 次数，让 CPU 和设备并行工作。

\begin{lstlisting}
// 注释（推进）：逐行追踪发起者、所有权和可见性；请求被接受不等于事务完成。
// 注释（边界）：以采样沿、状态位或 completion 为准，再允许消费者使用结果。
receive ring:
  CPU owns descriptor N
  CPU fills buffer address and length
  CPU marks descriptor ready for device
  CPU updates tail pointer or rings doorbell
  device writes packet into buffer
  device writes completion status
  CPU regains ownership and processes buffer
\end{lstlisting}

\begin{longtable}{L{0.18\linewidth}L{0.38\linewidth}L{0.34\linewidth}}
\toprule
环队列字段 & 硬件含义 & 关键问题 \\
\midrule
描述符地址 & 设备要读写的主存位置 & 地址是否经过 IOMMU 或物理地址转换 \\\tablerowrule
长度 & 本次 DMA 允许访问的字节数 & 设备不能越界写出缓冲区 \\\tablerowrule
所有权位 & 当前由 CPU 还是设备拥有 & 所有权不重叠，状态切换有屏障 \\\tablerowrule
头尾指针 & CPU 和设备约定的生产/消费位置 & 指针更新顺序和环绕规则清楚 \\\tablerowrule
completion & 设备报告完成、长度、错误码 & 中断到达不等于所有 cache 可见性自动正确 \\
\bottomrule
\end{longtable}

这个模型也说明驱动程序为什么必须理解硬件。驱动不是“调用设备 API”的普通软件，它要写 MMIO 寄存器、分配对齐缓冲区、建立 DMA 映射、维护 cache 可见性、处理中断、回收描述符和处理错误。

\topic{描述符环不动、指针动}
描述符环的关键是“环不动、指针动”：描述符数组在内存中的位置固定，CPU 只推进 SW 指针，设备只推进 HW 指针，走到末尾就回绕；两个指针的差值就是环中未完成事务的数量。doorbell 不携带数据，只是告诉设备“指针又动了”，设备即使错过一次 doorbell，靠指针差也能补读。反过来，某个描述符是否完成，不能凭 doorbell 或中断猜测，必须以设备写回描述符状态位（或 completion 队列项）为准；CPU 回收缓冲区前，还要确认这次写回对自己可见。

\begin{pdfdiagram}
  % 两行四格加回绕箭头，代替真圆环，保证正交走线
  \node[hw state, minimum width=2.0cm] (d0) at (1.3,2.6) {描述符 0};
  \node[hw state, minimum width=2.0cm] (d1) at (3.7,2.6) {描述符 1};
  \node[hw state, minimum width=2.0cm] (d2) at (6.1,2.6) {描述符 2};
  \node[hw state, minimum width=2.0cm] (d3) at (8.5,2.6) {描述符 3};
  \node[hw state, minimum width=2.0cm] (d7) at (1.3,1.2) {描述符 7};
  \node[hw state, minimum width=2.0cm] (d6) at (3.7,1.2) {描述符 6};
  \node[hw state, minimum width=2.0cm] (d5) at (6.1,1.2) {描述符 5};
  \node[hw state, minimum width=2.0cm] (d4) at (8.5,1.2) {描述符 4};
  \draw[hw bus, ->] (d0.east) -- (d1.west);
  \draw[hw bus, ->] (d1.east) -- (d2.west);
  \draw[hw bus, ->] (d2.east) -- (d3.west);
  \draw[hw bus, ->] (d3.south) -- (d4.north);
  \draw[hw bus, ->] (d4.west) -- (d5.east);
  \draw[hw bus, ->] (d5.west) -- (d6.east);
  \draw[hw bus, ->] (d6.west) -- (d7.east);
  \draw[hw bus, ->] (d7.north) -- (d0.south) node[midway, fill=white, inner sep=1pt, hw label] {回绕};
  % SW 指针（CPU 写到哪里）与 HW 指针（设备读到哪里）
  \draw[hw return, ->] (8.5,3.6) -- (d3.north);
  \node[hw label, above] at (8.5,3.6) {SW 指针：CPU 下一个写};
  \draw[hw return, ->] (3.7,3.6) -- (d1.north);
  \node[hw label, above] at (3.7,3.6) {HW 指针：设备正在读};
  % doorbell 写动作
  \node[hw control, minimum width=1.7cm] (bell) at (10.65,1.2) {doorbell};
  \draw[hw return, ->] (d3.east) -- (11.05,2.6) -- (11.05,1.51);
  \node[hw label, anchor=east] at (11.0,2.05) {写 doorbell};
\end{pdfdiagram}
\diagramnote{描述符环。设备按 HW 指针顺序消费，CPU 按 SW 指针顺序生产，描述符 7 之后回到 0；CPU 更新 SW 后写一次 doorbell 提示设备。完成与否以设备写回描述符状态位为准，doorbell 只是“有新描述符”的提示。}

\topic{驱动 bring-up 应分层验证}
调试设备时，最差的方法是直接运行完整驱动并等待它“能不能用”。更可靠的顺序是先验证配置空间或设备 ID，再验证 MMIO 映射，再验证单个寄存器读写，再验证中断，再验证一个最小 DMA 描述符，最后才跑持续队列。每一层都要有可观察的结果。

\begin{longtable}{L{0.18\linewidth}L{0.42\linewidth}L{0.3\linewidth}}
\toprule
bring-up 阶段 & 操作 & 预期结果 \\
\midrule
枚举/识别 & 读取设备 ID、版本或能力寄存器 & 地址映射正确，读值稳定可复现 \\\tablerowrule
复位 & 写复位位并轮询 ready/busy & 状态位按规格变化，不永久 busy \\\tablerowrule
MMIO 回读 & 写可回读配置寄存器再读回 & 字节序、位域和写掩码正确 \\\tablerowrule
中断 & 人工触发或配置定时事件 & 中断线/MSI 到达，清除顺序正确 \\\tablerowrule
单次 DMA & 一个描述符、一个缓冲区、一个 completion & 地址、长度、cache 可见性和错误码正确 \\\tablerowrule
持续队列 & 多描述符环和批量 completion & 所有权、环绕、背压和中断聚合正确 \\
\bottomrule
\end{longtable}

这张表也适用于教学板。即使没有 PCIe，只用一个 GPIO、定时器或简化 DMA 控制器，也应该按同样顺序验证：先确认地址译码，再确认寄存器语义，再确认事件通知，再确认数据搬运。

\topic{案例：一个接收包怎样进入内存}
以网卡接收一个数据包为例。CPU 先分配若干缓冲区，把每个缓冲区的地址和长度写入接收描述符环；随后保证描述符对设备可见，再写 MMIO 门铃或尾指针告诉设备可以接收。设备收到外部数据后，作为 DMA 发起者把数据写入某个缓冲区，再写 completion 或描述符状态，最后触发中断。CPU 处理中断，确认 completion，按一致性规则让缓冲区数据对 CPU 可见，然后把缓冲区交给上层软件。

\begin{longtable}{L{0.14\linewidth}L{0.42\linewidth}L{0.34\linewidth}}
\toprule
阶段 & 硬件事务 & 观察要点 \\
\midrule
准备缓冲区 & CPU 在 RAM 中分配对齐缓冲区 & 物理地址或 I/O 虚拟地址有效，长度足够 \\\tablerowrule
写描述符 & CPU STORE 地址、长度、所有权位 & 描述符 cache line 已按平台规则可见 \\\tablerowrule
通知设备 & CPU 写 MMIO tail/doorbell & MMIO 不被普通 cache 合并或延迟越界 \\\tablerowrule
设备 DMA 写 & 设备发起内存写事务 & IOMMU 权限、字节数、completion 顺序正确 \\\tablerowrule
投递完成 & 设备写状态并触发中断 & completion 可读，中断不会早于数据可见 \\\tablerowrule
CPU 回收 & 中断处理读取状态和缓冲区 & cache invalidate 或 coherent 规则完成 \\
\bottomrule
\end{longtable}

\begin{lstlisting}
// 注释（推进）：逐行追踪发起者、所有权和可见性；请求被接受不等于事务完成。
// 注释（边界）：以采样沿、状态位或 completion 为准，再允许消费者使用结果。
receive path:
  CPU: fill rx_desc[N] with buffer address and length
  CPU: barrier, then MMIO[RX_TAIL] = N
  DEV: reads rx_desc[N]
  DEV: writes packet bytes into buffer
  DEV: writes completion status
  DEV: raises interrupt
  CPU: acknowledges interrupt and reads completion
  CPU: synchronizes buffer visibility
  CPU: consumes packet
\end{lstlisting}

这条路径里至少有三次“完成”。第一，CPU 写 MMIO 门铃完成，只表示设备接口看到了通知；第二，设备 DMA 写缓冲区完成，表示数据已经到达内存系统的某个可见点；第三，中断处理完成，表示 CPU 已经确认状态并回收所有权。把这三者混在一起，是驱动和硬件协同时最常见的错误来源。

\topic{案例：发送路径为什么要区分数据和描述符}
发送数据包时方向相反。CPU 先把待发送数据写入缓冲区，再写发送描述符，最后通知设备。设备读描述符，再 DMA 读缓冲区，把数据发出。这里最关键的可见性顺序是：缓冲区内容必须先对设备可见，描述符必须描述正确地址和长度，门铃必须最后到达。否则设备可能发出旧数据、半写数据或错误长度的数据。

\begin{longtable}{L{0.18\linewidth}L{0.4\linewidth}L{0.32\linewidth}}
\toprule
发送对象 & 可见性要求 & 失败表现 \\
\midrule
数据缓冲区 & CPU 写入的数据先对设备可见 & 设备发送旧包或未初始化字节 \\\tablerowrule
描述符 & 地址、长度、标志位完整可见 & 设备读错缓冲区或长度 \\\tablerowrule
门铃寄存器 & 在数据和描述符之后对设备可见 & 设备提前启动，读到旧状态 \\\tablerowrule
completion & 设备确认已不再读取缓冲区 & CPU 过早重用缓冲区导致数据损坏 \\
\bottomrule
\end{longtable}

发送 completion 到达前，CPU 不应重用缓冲区。即使 MMIO 写门铃已经返回，设备可能还没有读描述符；即使设备已经读描述符，可能还没有读完数据；即使数据已经发出，completion 可能还在队列里等待 CPU 处理。因此所有权位和 completion 队列不是形式主义，而是让 CPU 和设备不同时修改同一块内存的硬件约定。

\topic{案例：块设备读请求的队列化}
磁盘、NVMe、SD 控制器或简化块设备也可以用同一模型理解。CPU 准备命令描述符：逻辑块地址、目标缓冲区、长度、读写方向；设备读取命令后，在内部介质或控制器中完成实际读写，再通过 DMA 把数据写到主存，最后投递 completion。区别只在于设备内部工作可能更慢、队列更深、错误码更复杂。

\begin{longtable}{L{0.18\linewidth}L{0.38\linewidth}L{0.34\linewidth}}
\toprule
块请求字段 & 含义 & 硬件问题 \\
\midrule
LBA 或块号 & 设备内部介质位置 & 越界、坏块、介质错误如何报告 \\\tablerowrule
缓冲区地址 & DMA 写入或读取的主存位置 & 对齐、IOMMU、cache 可见性 \\\tablerowrule
长度 & 传输字节数或块数 & 不能越过缓冲区和描述符许可范围 \\\tablerowrule
方向 & 读设备到内存，或写内存到设备 & DMA 方向决定 cache 维护方式 \\\tablerowrule
命令 ID & completion 对应哪个请求 & 队列乱序完成时仍能匹配请求 \\\tablerowrule
状态码 & 成功、超时、校验错误、权限错误 & 驱动恢复或上报错误 \\
\bottomrule
\end{longtable}

这个案例能把“设备很慢”讲得更具体。CPU 发出块读请求后，并不会逐字等待数据从设备口读出。它可以切换去做别的工作；设备和 DMA 在后台推进；中断或轮询 completion 时，CPU 再取回结果。硬件性能来自异步队列和批量搬运，而不是让单条 LOAD 直接等磁盘或网络。

\topic{案例：一个未映射访问如何成为可诊断错误}
再看失败路径。CPU 执行 \code{MOV R1,[0xDEAD0000]}，地址译码没有任何目标响应。没有错误约定时，CPU 可能读到悬空总线值，或永久等待 ready。可靠平台应规定超时、总线错误或异常入口。若有 MMU，还可能在页表阶段就发现该地址未映射，不会走到外部总线。

\begin{longtable}{L{0.2\linewidth}L{0.38\linewidth}L{0.32\linewidth}}
\toprule
失败位置 & 可观测现象 & 正确处理 \\
\midrule
页表不存在 & 地址转换阶段触发页错误 & 保存故障地址和错误码，进入页错误入口 \\\tablerowrule
权限不允许 & 转换存在但读写/特权不满足 & 触发保护异常，不发起外部访问 \\\tablerowrule
地址译码无目标 & 总线事务无人响应 & 超时或 bus error，而不是无限等待 \\\tablerowrule
目标返回错误 & 设备或内存控制器报告失败 & 错误状态进入异常或驱动恢复路径 \\\tablerowrule
ECC 校验失败 & 数据返回但校验错误 & 更正、重读、上报或机器检查 \\
\bottomrule
\end{longtable}

现象“程序卡住”可能来自页表、总线译码、设备 ready、时钟、复位或错误入口任何一层。把失败位置写成层级，就能决定先看页表项、地址线、片选、ready、错误状态，还是异常向量。

\begin{classicnote}
读任何硬件平台，都可以用同一张事务表约束理解：发起者、目标、地址、方向、数据宽度、握手、完成、错误、副作用和可观测状态。早期 8086 总线、微控制器片内外设、PCIe 设备和现代 SoC 互连的规模不同，但这十项问题不变。
\end{classicnote}

\topic{专题：DRAM 控制器为什么要排队和调度}
SRAM 可以近似看成“地址稳定后读出数据”的存储器，DRAM 则必须先打开行、读写列、预充电、刷新。主存控制器面对的不是单个 LOAD，而是一串来自 CPU、DMA、显示控制器和其他主设备的请求。它要在正确性约束内排序请求，尽量利用行命中和 bank 并行，同时保证刷新不丢失数据。

\begin{longtable}{L{0.18\linewidth}L{0.42\linewidth}L{0.3\linewidth}}
\toprule
DRAM 概念 & 含义 & 对系统的影响 \\
\midrule
row activate & 把某一行搬到行缓冲中 & 同一行后续列访问较快 \\\tablerowrule
row conflict & 下一个请求落在不同的行 & 需要预充电再打开新行，延迟增加 \\\tablerowrule
bank & 多组可相对独立工作的阵列 & 交错访问可隐藏部分等待 \\\tablerowrule
refresh & 周期性恢复电容电荷 & 刷新窗口内部分请求会被延后 \\\tablerowrule
ECC & 用冗余位检测或纠正错误 & 增加位宽和延迟，但提高可靠性 \\
\bottomrule
\end{longtable}

因此，主存延迟不是一个固定常数。一次 cache miss 到达内存控制器后，可能遇到行命中、行冲突、刷新、写队列排空或高优先级 DMA。系统软件看到的是“同样的 LOAD 有时快有时慢”，硬件工程师看到的是请求队列、调度策略和时序参数。

\topic{专题：cache 设计参数怎样改变命中率和时延}
cache 的基本问题是把大而慢的主存切成较小的行，并保存最近使用的数据。地址会被拆成 tag、index、offset：offset 选择 cache 行内字节，index 选择组，tag 判断这一组里哪一路是目标行。不同组织改变冲突概率、比较器数量、替换策略和命中时延。

\begin{longtable}{L{0.18\linewidth}L{0.36\linewidth}L{0.36\linewidth}}
\toprule
组织 & 优点 & 代价或风险 \\
\midrule
直接映射 & 查找简单、速度快、面积小 & 两个常用地址若映到同一行会反复冲突 \\\tablerowrule
组相联 & 多路候选降低冲突 miss & 多个 tag 比较器和路选择 mux 增加延迟 \\\tablerowrule
全相联 & 任意行可放任意位置，冲突最少 & 硬件代价高，通常用于 TLB 或小 cache \\\tablerowrule
写穿 & 内存总是较新，恢复简单 & 写流量大，常需写缓冲 \\\tablerowrule
写回 & 多次写合并到 cache 行，带宽效率高 & dirty 行替换时必须写回，失电/一致性更复杂 \\
\bottomrule
\end{longtable}

评价 cache 不能只看“第二次访问变快”。还要看替换、脏行写回、字节写使能、未对齐访问、MMIO 不可 cache、DMA 可见性和异常路径。例如设备把数据 DMA 到主存后，CPU 若继续读旧 cache 行，就会看到过期数据；CPU 写好发送缓冲区后，若脏行没有写回，设备也可能读到旧内容。这就是为什么驱动要有 cache clean/invalidate 或一致性互连。

\topic{专题：MMIO 寄存器位语义要像数据手册一样写}
MMIO 寄存器不是普通变量。每一位都可能有副作用：读清、写 1 清、只读、保留位必须写 0、写后自动启动、硬件置位软件清除。

\begin{longtable}{L{0.18\linewidth}L{0.36\linewidth}L{0.36\linewidth}}
\toprule
位语义 & 例子 & 软件访问要求 \\
\midrule
RO & 输入电平、设备 ID & 写入无效或保留，不应读改写破坏其他位 \\\tablerowrule
WO & FIFO 写口、启动命令 & 读值可能无意义，不能依赖回读 \\\tablerowrule
RW & 配置模式、使能位 & 修改前要明确复位默认值 \\\tablerowrule
W1C & 中断状态位 & 写 1 清除对应事件，写 0 保持 \\\tablerowrule
RC & 读状态自动清除 & 调试打印也可能改变设备状态 \\\tablerowrule
reserved & 未来扩展或工厂测试位 & 按手册写固定值，通常为 0 \\
\bottomrule
\end{longtable}

一个严谨的 GPIO 寄存器描述至少要写清：方向寄存器决定引脚驱动还是输入；输出寄存器只影响输出模式；输入寄存器来自同步后的外部电平；中断状态由边沿检测置位；清除中断采用 W1C；保留位读值未定义且写 0。这样驱动代码才能避免“读状态时把事件清掉”“读改写时误清中断”“把保留位写成 1 后芯片行为改变”等问题。

\topic{专题：总线错误要能定位到层级}
当一次访问失败时，工程上最重要的问题不是“错了”，而是错在地址转换、权限、互连、目标设备还是数据校验。若每层都只把错误折叠成一个笼统异常，调试会非常困难。更好的设计是保留故障地址、访问方向、主设备 ID、错误来源和目标返回码。

\begin{longtable}{L{0.2\linewidth}L{0.42\linewidth}L{0.28\linewidth}}
\toprule
记录字段 & 作用 & 调试价值 \\
\midrule
fault address & 失败访问的地址 & 区分空指针、越界和 MMIO 映射错误 \\\tablerowrule
access type & 读、写、取指、原子操作 & 判断权限和副作用风险 \\\tablerowrule
master ID & CPU、DMA、GPU 或调试器 & 找到真正发起者 \\\tablerowrule
decode result & 命中哪个目标或未命中 & 定位地址地图错误 \\\tablerowrule
response code & OKAY、SLVERR、DECERR、超时等 & 区分目标报错和无人响应 \\
\bottomrule
\end{longtable}

这也是为什么后续平台 bring-up 日志要记录总线层的事务信息。单看 C 代码中的一行 \code{*ptr} 无法知道失败来自页表、cache、IOMMU、互连还是设备时钟未开；完整错误记录能把问题缩小到可测量的硬件边界。

\topic{专题：内存层次要同时看容量、时延、带宽和所有权}
计算机硬件不是只有 CPU。CPU 的每一次有用工作，几乎都要和某一级存储或 I/O 交换信息。寄存器最快但数量极少；L1/L2/L3 cache 容量逐级增大、时延逐级上升；DRAM 容量大但需要控制器、队列和刷新；SSD/硬盘容量更大，却只能通过控制器、队列、DMA 和中断间接参与执行。把这些对象都称为“存储”会掩盖关键差异：有些层按字节随机访问，有些层按 cache line 交换，有些层按页或块传输，有些层只能通过命令队列返回 completion。

\begin{pdfdiagram}[x=1cm,y=1cm]
  \node[hw state, minimum width=1.55cm] (reg) at (0,0) {寄存器};
  \node[hw compute, minimum width=1.55cm] (l1) at (1.8,0) {L1};
  \node[hw compute, minimum width=1.55cm] (l2) at (3.6,0) {L2/L3};
  \node[hw memory, minimum width=1.75cm] (dram) at (5.65,0) {DRAM};
  \node[hw control, minimum width=1.75cm] (ssd) at (7.9,0) {SSD/硬盘};
  \node[hw lane, minimum width=8.8cm, text width=8.4cm] at (3.95,-1.35) {越靠左，访问粒度越小、时延越低、容量越小；越靠右，容量越大，但访问要经过控制器、队列、DMA、错误码和完成通知。};
  \draw[hw bus] (reg) -- node[above,hw label]{操作数} (l1);
  \draw[hw bus] (l1) -- node[above,hw label]{cache line} (l2);
  \draw[hw bus] (l2) -- node[above,hw label]{回填/写回} (dram);
  \draw[hw bus] (dram) -- node[above,hw label]{DMA 缓冲区} (ssd);
\end{pdfdiagram}
\diagramnote{内存层次不是一条“更大的数组”，而是一组访问粒度和完成语义不同的层。}

\begin{longtable}{L{0.16\linewidth}L{0.22\linewidth}L{0.28\linewidth}L{0.24\linewidth}}
\toprule
层次 & 常见粒度 & 谁直接访问 & 主要硬件问题 \\
\midrule
寄存器 & word 或向量 lane & 执行单元 & 端口数、旁路、写回冲突 \\\tablerowrule
L1 cache & cache line & 当前核心 load/store 单元 & 命中时延、虚实地址、别名和写策略 \\\tablerowrule
L2/L3 cache & cache line & 一个核心或多个核心 & 容量、共享、替换、一致性 \\\tablerowrule
DRAM & burst/cache line & 内存控制器 & bank、行命中、刷新、ECC、调度 \\\tablerowrule
SSD/硬盘 & 块、页或扇区 & 设备控制器和 DMA & 命令队列、坏块、磨损、错误恢复 \\\tablerowrule
网络/显卡等设备 & 描述符、包、帧、buffer & DMA 引擎和设备内部队列 & 所有权、completion、IOMMU、cache 可见性 \\
\bottomrule
\end{longtable}

从程序角度看，\code{load x} 只是取一个变量；从硬件角度看，它可能命中寄存器旁路、命中 L1、命中 L2、命中共享 L3、访问 DRAM、触发页表遍历、发生缺页，甚至等待块设备把页面读回主存。《深入理解计算机系统》（CSAPP）强调的 locality 就是让常见访问停留在左侧短路径。

\topic{专题：虚拟地址别名与同源不同名的 cache 问题}
上一专题的层次表把“别名”列为 L1 cache 的关键问题之一，这里把它展开。TLB 的存在意味着同一个物理页可以被多个虚拟页映射：共享库、进程间共享内存、同一页在用户态和内核态各有别名，都是日常情形。这带来一个只在“用虚拟地址索引 cache”时才出现的问题：两个不同的虚拟地址指向同一物理行，若 cache 按虚拟地址的低位选槽，而两个别名的虚拟页号不同，index 就可能不同，同一物理行会被存进两个不同的槽。之后程序经别名一写入，经别名二读到的却是另一个槽里的旧副本——同一数据在 cache 里有了两份不一致的拷贝。

看一个小例子。页大小 4 KiB，cache 为 32 KiB 直接映射、64 B 行：offset 为 6 bit，行数 $=32\ \text{KiB}\div 64\ \text{B}=512$，index 取地址的 bit 14..6 共 9 bit。页内偏移只有 12 bit（bit 11..0），因此 index 的高 3 位（bit 12..14）来自虚拟页号——别名地址的差异正好在这几位里。设虚拟地址 \code{0x1040} 和 \code{0x2040} 都映射到物理地址 \code{0x35040}：前者的 index 是 65，后者的 index 是 129，同一物理行会被分别存进槽 65 和槽 129。经 \code{0x1040} 写入后，槽 65 是新值，槽 129 仍是旧值，两次读到同一物理地址却得到不同数据。

两种硬件解法基于同一条思路：让同一物理行的所有别名映到同一个槽，再由物理 tag 确认身份。第一种是物理索引（PIPT）：index 直接取自物理地址，\code{0x1040} 和 \code{0x2040} 经 TLB 翻译后都是 \code{0x35040}，index 都是 321，必然同一个槽，tag 也比较物理地址，别名问题根本不会出现；代价是要先等 TLB 翻译完成才能开始查 cache，命中路径变长。第二种是限制 index 位数（VIPT 的常见约定）：把 cache 设计成 index 加 offset 不超过页内偏移的 12 bit，例如 4 KiB 直接映射，或 32 KiB、8 路组相联——组数 $=32\ \text{KiB}\div 64\ \text{B}\div 8=64$，index 只需 6 bit，index 加 offset 正好 12 bit。此时 index 全部来自页内偏移，而任何别名的页内偏移都与物理地址相同：上例中两个别名与物理地址的 bit 11..6 都等于 1，必然映到同一组。查找与 TLB 翻译并行进行，拿到物理页号后再与组内各路的物理 tag 比较，第二次访问命中同一个行，两份副本无从产生。现代处理器的 L1 大多采用这种“虚拟索引、物理 tag、index 不越出页内偏移”的安排，既保留并行查找的速度，又让别名退化为普通命中。

\begin{pdfdiagram}
  % 左：未限制 index 位数（32 KiB 直接映射，index 取 bit 14..6）
  \node[hw label] at (1.9,2.75) {未限制：按虚拟 index 映到两个槽};
  \node[hw state, minimum width=1.6cm] (va1) at (-0.3,0.9) {VA \code{0x1040}};
  \node[hw state, minimum width=1.6cm] (va2) at (-0.3,-0.9) {VA \code{0x2040}};
  \node[hw compute, minimum width=1.0cm] (tlb) at (1.8,0) {TLB};
  \node[hw state, minimum width=1.9cm] (pa) at (3.95,0) {PA \code{0x35040}};
  \node[hw memory, minimum width=2.0cm] (s65) at (3.95,1.9) {槽 65：新值};
  \node[hw memory, minimum width=2.0cm] (s129) at (3.95,-1.9) {槽 129：旧值};
  \draw[hw bus] (va1.north) -- (-0.3,1.9) -- node[pos=0.7,above=1pt,hw label,fill=white,inner sep=1pt]{index 65} (s65.west);
  \draw[hw bus] (va2.south) -- (-0.3,-1.9) -- node[pos=0.7,above=1pt,hw label,fill=white,inner sep=1pt]{index 129} (s129.west);
  \draw[hw bus] (va1.east) -| (tlb.north);
  \draw[hw bus] (va2.east) -| (tlb.south);
  \draw[hw bus] (tlb) -- node[above=1pt,hw label,fill=white,inner sep=1pt]{同一 PA} (pa);
  \draw[hw return] (pa.north) -- node[right=1pt,hw label]{tag} (s65.south);
  \draw[hw return] (pa.south) -- node[right=1pt,hw label]{tag} (s129.north);
  \node[hw label] at (1.9,-2.75) {同一物理行出现两份副本};
  % 右：限制 index 位数（index+offset 不越出页内偏移 12 bit）
  \node[hw label] at (7.3,2.75) {限制 index 位数：别名同组};
  \node[hw state, minimum width=1.9cm] (var) at (6.2,0.9) {VA \code{0x1040}\\VA \code{0x2040}};
  \node[hw compute, minimum width=1.0cm] (tlb2) at (8.3,0.9) {TLB};
  \node[hw memory, minimum width=2.0cm] (grp) at (8.3,-0.9) {组 1：唯一副本};
  \draw[hw bus] (var) -- (tlb2);
  \draw[hw bus] (var.south) -- node[pos=0.4,right=1pt,hw label]{index = bit 11..6} (6.2,-0.9) -- (grp.west);
  \draw[hw return] (tlb2.south) -- node[right=1pt,hw label]{PA \code{0x35040}\\物理 tag} (grp.north);
  \node[hw label] at (7.3,-2.75) {index 全来自页内偏移};
\end{pdfdiagram}
\diagramnote{VIPT 别名问题与限制 index 位数的对照。左：32 KiB 直接映射 cache 的 index 取 bit 14..6，\code{0x1040} 与 \code{0x2040} 的 index 分别是 65 与 129，同一物理行 \code{0x35040} 被存进两个槽，经 \code{0x1040} 写入后槽 65 是新值、槽 129 仍是旧值。右：限制 index 加 offset 不越出页内偏移 12 bit（如 32 KiB、8 路组相联，index 为 bit 11..6），两个别名的 index 都来自页内偏移，必然映到同一组，再由物理 tag 确认身份，双副本无从产生。}

\topic{专题：cache 与 TLB 是两套不同的近路}
cache 缓存数据或指令内容，TLB 缓存虚拟页到物理页的地址转换。二者经常同时出现在一次 LOAD 路径上，但回答的问题不同：TLB 问“这个虚拟地址能否翻译成哪个物理地址”；cache 问“这个物理位置的内容是否已经在近处”。TLB miss 不等于 cache miss，页错误也不等于总线错误。

\begin{pdfdiagram}
  \node[hw state, minimum width=1.85cm] (va) at (0,0) {虚拟地址};
  \node[hw compute, minimum width=1.85cm] (tlb) at (2.45,0) {TLB};
  \node[hw state, minimum width=1.85cm] (pa) at (4.9,0) {物理地址};
  \node[hw compute, minimum width=1.85cm] (cache) at (7.35,0) {cache};
  \node[hw memory, minimum width=1.65cm] (data) at (9.65,0) {数据};
  \node[hw control, minimum width=2.0cm] (ptw) at (2.45,-1.35) {页表遍历};
  \node[hw memory, minimum width=2.0cm] (mem) at (7.35,-1.35) {DRAM 回填};
  \draw[hw bus] (va) -- (tlb);
  \draw[hw bus] (tlb) -- node[above,hw label]{命中} (pa);
  \draw[hw bus] (pa) -- (cache);
  \draw[hw bus] (cache) -- node[above,hw label]{命中} (data);
  \draw[hw danger] (tlb) -- node[right,hw label]{miss} (ptw);
  \draw[hw danger] (cache) -- node[right,hw label]{miss} (mem);
\end{pdfdiagram}
\diagramnote{TLB miss 发生在地址转换阶段，cache miss 发生在数据回填阶段。二者都可能让 LOAD 变慢，但失败位置不同。}

\begin{longtable}{L{0.2\linewidth}L{0.34\linewidth}L{0.36\linewidth}}
\toprule
事件 & 真实含义 & 后续动作 \\
\midrule
TLB hit + cache hit & 地址转换和数据都在近处 & LOAD 快速返回 \\\tablerowrule
TLB miss + 页表命中 & 转换不在 TLB，但页表有效 & 硬件或软件页表遍历后填 TLB \\\tablerowrule
页不存在 & 页表没有有效映射 & 进入页错误异常，可能由 OS 分配或从磁盘读页 \\\tablerowrule
cache miss & 地址有效，但数据不在当前 cache & 从下一级 cache 或 DRAM 回填 cache line \\\tablerowrule
权限错误 & 转换存在但访问类型不允许 & 报保护异常，不能提交数据 \\
\bottomrule
\end{longtable}

\topic{专题：SSD 和硬盘不是慢内存，而是块设备}
硬盘和 SSD 都不能像 SRAM 那样把地址线接上就读一个字节。CPU 要通过控制器提交命令：读哪个逻辑块、读多长、放到哪个主存缓冲区、完成后写哪个 completion。硬盘的内部约束来自机械寻道和旋转；SSD 的内部约束来自 NAND flash 页读写、块擦除、FTL（闪存转换层，flash translation layer）映射、磨损均衡、垃圾回收和掉电保护。它们对 CPU 暴露的共同形态，是队列化块设备。

\begin{longtable}{L{0.18\linewidth}L{0.34\linewidth}L{0.38\linewidth}}
\toprule
对象 & 内部结构重点 & 对系统可见的边界 \\
\midrule
机械硬盘 & 盘片、磁头、磁道、扇区、缓存 & 随机访问受寻道和旋转影响，顺序访问更友好 \\\tablerowrule
SATA SSD & NAND flash、FTL、通道、擦除块 & 写放大、垃圾回收、坏块管理、掉电保护 \\\tablerowrule
NVMe SSD & PCIe、多队列、doorbell、completion queue & 并发队列深度高，CPU 通过 MMIO 和 DMA 协作 \\\tablerowrule
SD/eMMC & 控制器隐藏 flash 细节，命令集较简单 & 嵌入式启动和存储，延迟波动明显 \\
\bottomrule
\end{longtable}

理解 SSD 时要避免两个误解。第一，SSD 没有机械寻道，不等于每次访问都一样快；内部垃圾回收、擦除、映射表缓存和热控都会让尾延迟上升。第二，逻辑块地址连续，不等于物理 NAND 页连续；FTL 会把逻辑写重映射到新位置，再在后台回收旧块。因此块设备的正确抽象不是“巨大数组”，而是“命令队列 + DMA 缓冲区 + completion + 错误恢复”。

\begin{pdfdiagram}
  \node[hw state, minimum width=1.55cm] (cpu) at (0,0) {CPU};
  \node[hw control, minimum width=2.05cm] (sq) at (2.55,0.72) {提交队列};
  \node[hw control, minimum width=1.65cm] (bell) at (2.55,-0.72) {门铃};
  \node[hw control, minimum width=2.0cm] (ctrl) at (5.25,0) {NVMe 控制器};
  \node[hw memory, minimum width=1.85cm] (nand) at (7.8,0.72) {NAND/FTL};
  \node[hw state, minimum width=2.05cm] (cq) at (7.8,-0.72) {完成队列};
  \node[hw label] at (1.25,1.15) {1 写命令};
  \node[hw label] at (1.25,-1.12) {2 敲门铃};
  \node[hw label] at (6.55,-1.18) {3 投递完成};
  \draw[hw bus] (cpu) -- (sq);
  \draw[hw bus] (cpu) -- (bell);
  \draw[hw bus] (sq) -- (ctrl);
  \draw[hw bus] (bell) -- (ctrl);
  \draw[hw bus] (ctrl) -- (nand);
  \draw[hw return] (ctrl) -- (cq);
  \draw[hw return] (cq.west) |- (cpu.south);
\end{pdfdiagram}
\diagramnote{NVMe 的基本形态：CPU 写队列和门铃，设备 DMA 读写主存，completion 才说明请求完成。}

\topic{专题：读硬件结构时要区分数据路径和控制路径}
CPU、cache、DRAM、SSD、网卡和显卡都可以用同一个提问框架阅读：数据从哪里到哪里，控制由谁发起，完成由谁报告，错误由谁保存。低质量描述常把这些混成“设备把数据给 CPU”。更好的描述应写出两条路径：数据可能经 DMA 在设备和主存之间移动；控制通常由 CPU 写 MMIO、设备写 completion、CPU 读状态或处理中断。

\begin{longtable}{L{0.18\linewidth}L{0.36\linewidth}L{0.36\linewidth}}
\toprule
路径 & 典型方向 & 必须说明 \\
\midrule
控制路径 & CPU 到 MMIO/doorbell，设备到状态/中断 & 命令何时被接收，完成如何通知 \\\tablerowrule
数据路径 & 设备 DMA 与主存缓冲区之间 & 缓冲区地址、长度、cache 可见性、权限 \\\tablerowrule
错误路径 & 设备或互连到错误码/异常 & 超时、校验、越界、权限和重试策略 \\\tablerowrule
恢复路径 & 驱动或固件到 reset/retry/rebuild & 哪些状态保留，哪些队列丢弃或重放 \\
\bottomrule
\end{longtable}
































































































\section{四、存储层次总览}

把本章所有器件放在同一张表中，可以看到物理结构如何决定系统参数：

\begin{longtable}{L{0.11\linewidth}L{0.13\linewidth}L{0.12\linewidth}L{0.12\linewidth}L{0.12\linewidth}L{0.14\linewidth}L{0.13\linewidth}}
\toprule
器件 & 物理机制 & 容量 & 读延迟 & 写延迟 & 保持 & 失效模式 \\
\midrule
寄存器 & 触发器 & KB & $<$1 ns & $<$1 ns & 通电 & 断电丢失 \\\tablerowrule
SRAM (L1/L2) & 6T 锁存 & 32K--16M & 1--4 ns & 1--4 ns & 通电 & 断电丢失；SNM 不足→读破坏/写失败；$\alpha$/中子→软错误 \\\tablerowrule
DRAM & 1T1C 电容 & 8--128 GB & 50--100 ns & 50--100 ns & 需刷新 & 漏电→数据丢失；刷新失败 \\\tablerowrule
NAND SSD & 浮栅/电荷俘获 & 256G--8T & 25--100 $\mu$s & 200 $\mu$s--2 ms & 10 年（新片常温；寿命末期消费级 1 年 @ 30°C） & P/E 耗尽→坏块；读干扰 \\\tablerowrule
HDD & 磁畴方向 & 1--24 TB & 4--12 ms & 4--12 ms & 数十年 & 磁头碰撞；伺服丢失；介质退化 \\
\bottomrule
\end{longtable}

每一层的访问粒度、写入代价和失效模式都由其物理结构决定，不能通过软件抽象消除。理解这些结构，才能理解为什么 cache miss 代价是 100 ns 而不是 1 ns，为什么 SSD 随机写比顺序写慢，为什么 HDD 随机 IOPS 只有 100 而顺序带宽有 200 MB/s。

\section*{章末练习}

\begin{chapterexercise}{SRAM 读写周期}
画出地址、CS、OE、WE、数据线和采样窗口
\exercisecheck{预期结果}{能指出谁驱动数据线，谁采样}
\end{chapterexercise}

\begin{chapterexercise}{地址映射}
设计 64 KiB ROM/RAM/GPIO/定时器映射
\exercisecheck{预期结果}{任意地址最多一个片选有效}
\end{chapterexercise}

\begin{chapterexercise}{等待状态}
为慢 ROM 和慢 I/O 设计 READY/WAIT 规则
\exercisecheck{预期结果}{等待期间地址和方向稳定，有超时策略}
\end{chapterexercise}

\begin{chapterexercise}{MMIO 寄存器}
设计 GPIO 的 \code{DIR/OUT/IN/STAT} 语义
\exercisecheck{预期结果}{外部输入经过同步，状态清除规则明确}
\end{chapterexercise}

\begin{chapterexercise}{中断顺序}
写出设备 ready 到 CPU 服务程序返回的 trace
\exercisecheck{预期结果}{不丢事件，不重复清旧事件}
\end{chapterexercise}

\begin{chapterexercise}{DMA 顺序}
写出描述符、缓冲区、completion 和 cache 可见性边界
\exercisecheck{预期结果}{CPU 和设备所有权不重叠}
\end{chapterexercise}

\begin{chapterexercise}{valid/ready}
用 valid/ready 描述一次读请求和读响应
\exercisecheck{预期结果}{能说明背压和 completion 的作用}
\end{chapterexercise}

\begin{chapterexercise}{cache 行}
用 tag/index/offset 解释一次 cache 命中和一次 miss
\exercisecheck{预期结果}{能说明 cache line、valid、dirty 和替换边界}
\end{chapterexercise}

\begin{chapterexercise}{MMIO 属性}
说明为什么设备寄存器通常不能按普通 cacheable 内存处理
\exercisecheck{预期结果}{能指出副作用、顺序和可见性要求}
\end{chapterexercise}

\begin{chapterexercise}{未映射访问}
规定未映射地址的超时、错误或固定返回值
\exercisecheck{预期结果}{CPU 不会永久等待未知目标}
\end{chapterexercise}

\begin{chapterexercise}{存储层次图}
画出寄存器、L1/L2/L3、DRAM、SSD/硬盘的数据粒度和完成语义
\exercisecheck{预期结果}{能区分 cache line、页、块和 completion}
\end{chapterexercise}

\begin{chapterexercise}{TLB/cache 区分}
写出一次 LOAD 中 TLB hit/miss、页错误、cache miss 的不同路径
\exercisecheck{预期结果}{不把地址转换错误和数据回填混为一谈}
\end{chapterexercise}

\begin{chapterexercise}{块设备队列}
为 NVMe 或简化 SSD 写出 submission、doorbell、DMA、completion 顺序
\exercisecheck{预期结果}{能说明 CPU 写门铃不等于 I/O 完成}
\end{chapterexercise}

\begin{chapterexercise}{DRAM 时序计算}
给定 tRCD、tCL、时钟周期和突发长度，估算一次随机读时延，并与行命中情形对比
\exercisecheck{预期结果}{三段时延拆分正确，量级为几十 ns}
\end{chapterexercise}

\begin{chapterexercise}{UART 帧分析}
给定 RX 波形和标称波特率，数出起始位、数据位、停止位，还原字节并核对波特率误差
\exercisecheck{预期结果}{帧结构和误差判断都有依据}
\end{chapterexercise}

\begin{chapterexercise}{组相联查找}
用本章的四次访问序列，分别在直接映射和 2 路组相联下走一遍
\exercisecheck{预期结果}{能指出哪几次由 miss 变 hit 及原因}
\end{chapterexercise}

\begin{chapterexercise}{描述符环顺序}
写出 CPU 提交两个描述符到设备完成的 SW/HW 指针、doorbell 与完成写回顺序
\exercisecheck{预期结果}{doorbell 不当作完成，写回可见性明确}
\end{chapterexercise}

\begin{chapterexercise}{bank 交错拍数}
4 个 bank 按 cache line 轮转交错，单 bank 一次访问占用 4 拍，bank 间每拍可启动一个新请求；计算连续读 8 条相邻 cache line 的总拍数，并与单 bank 串行对比
\exercisecheck{预期结果}{交错 11 拍、单 bank 32 拍，能写出每行的启动与完成拍}
\end{chapterexercise}

\begin{chapterexercise}{替换策略}
2 路组相联、LRU，某组初始为空，依次访问 A、B、A、C、A、B；逐步写出命中/失效和组内容，并与同 index 直接映射对比
\exercisecheck{预期结果}{2 路共 3 次 miss，直接映射 6 次全 miss}
\end{chapterexercise}

\begin{chapterexercise}{预取演算}
顺序读 16 条 cache line，每次 miss 延迟 30 拍、每行处理 8 拍；计算无预取与 next-line 预取的总拍数，并求能完全隐藏预取的最小处理节拍
\exercisecheck{预期结果}{608 拍对 488 拍，处理节拍不小于 30 才能完全隐藏}
\end{chapterexercise}

\begin{chapterexercise}{I2C 事务序列}
写出主机从 7 位地址 \code{0x50} 的从机寄存器 \code{0x10} 读 2 字节的完整序列：START、地址字节、读写位、ACK/NACK、repeated START、STOP
\exercisecheck{预期结果}{\code{0xA0}/\code{0xA1} 正确，最后一字节 NACK，repeated START 位置正确}
\end{chapterexercise}

\begin{chapterexercise}{仲裁对照}
用一句话分别说明集中式仲裁器、菊花链、轮询计数的线数、优先级和公平性，并指出片上互连中仲裁的位置
\exercisecheck{预期结果}{三种做法的取舍清楚，仲裁做进交换节点}
\end{chapterexercise}

\begin{chapterexercise}{MESI 走查}
两核初始均为 I，按“A 读 X、B 读 X、A 写 X、B 读 X”写出每步两核状态和互连消息
\exercisecheck{预期结果}{每步最新值位置唯一，失效与写回/转发正确}
\end{chapterexercise}

\section*{参考解答与要点}

\begin{chapteranswer}{SRAM 读写周期}
读时地址和片选稳定，OE 有效、WE 无效，由 SRAM 驱动数据；写时 OE 无效，由发起者驱动数据，WE 在地址和数据稳定后有效。
\answercheck{必须保留的要点}{输出使能互斥，写窗口满足建立/保持时间。}
\end{chapteranswer}

\begin{chapteranswer}{地址映射}
例如低半区 ROM，高半区中一块 RAM，最高 4 KiB 为 I/O；用高位地址产生互斥片选，未映射区返回错误或超时。
\answercheck{必须保留的要点}{片选互斥和未映射行为必须写入规格。}
\end{chapteranswer}

\begin{chapteranswer}{等待状态}
慢目标在 ready 之前保持事务未完成，控制器保持地址、片选和方向；超过最大等待则返回错误。
\answercheck{必须保留的要点}{等待状态是硬件事务延长。}
\end{chapteranswer}

\begin{chapteranswer}{MMIO 寄存器}
GPIO 输出写入锁存器，输入来自同步后的外部引脚，状态位说明 ready/error/irq，清除规则避免新旧事件混淆。
\answercheck{必须保留的要点}{MMIO 有副作用，不是普通变量。}
\end{chapteranswer}

\begin{chapteranswer}{中断顺序}
设备锁存事件并请求中断，中断控制器投递，CPU 进入入口，服务程序读状态、处理数据、按设备规则清事件，再发送 EOI 或恢复中断。
\answercheck{必须保留的要点}{清除顺序决定是否丢事件。}
\end{chapteranswer}

\begin{chapteranswer}{DMA 顺序}
CPU 准备描述符和缓冲区并保证对设备可见，设备搬运数据并投递 completion，CPU 在 completion 后按一致性规则读取缓冲区。
\answercheck{必须保留的要点}{描述符可见、数据可见和缓冲区所有权必须分开。}
\end{chapteranswer}

\begin{chapteranswer}{valid/ready}
只有 valid 和 ready 同时为 1 才传输；ready 为 0 时发送方保持 payload；读请求和读响应可隔多个周期。
\answercheck{必须保留的要点}{背压不等于错误，completion 才是完成标志。}
\end{chapteranswer}

\begin{chapteranswer}{cache 行}
先算 offset=log2（行大小）、index=log2（行数）、tag=剩余位；命中要求槽 valid 且 tag 相等；miss 时取整行填入并更新 tag/valid，dirty 位决定替换时是否要先写回下一级。
\answercheck{必须保留的要点}{三个字段的分工和 valid/dirty 语义不能丢。}
\end{chapteranswer}

\begin{chapteranswer}{MMIO 属性}
设备寄存器有副作用：写是命令、读可能清状态；若被 cache 缓冲或合并，设备可能永远看不到命令，CPU 也可能一直读到旧状态，所以设备内存要映射为不可缓存、不可合并、顺序保持。
\answercheck{必须保留的要点}{副作用、顺序、可见性三条都要答到。}
\end{chapteranswer}

\begin{chapteranswer}{未映射访问}
未映射区域可以返回固定值、产生错误或等待到超时，但规则必须固定且可观察；不能让总线悬空成为随机值。
\answercheck{必须保留的要点}{错误路径也是硬件规格。}
\end{chapteranswer}

\begin{chapteranswer}{存储层次图}
寄存器（字，当拍完成）、L1/L2/L3（cache line，命中即返回）、DRAM（行/页，控制器调度与刷新）、SSD/硬盘（块，I/O 请求和 completion）；越往下粒度越大，完成语义越依赖显式通知而不是当拍返回。
\answercheck{必须保留的要点}{粒度与完成语义要对应，不能只画容量和速度。}
\end{chapteranswer}

\begin{chapteranswer}{TLB/cache 区分}
TLB miss 发生在地址转换阶段：硬件查页表，缺页则走页错误异常；cache miss 发生在拿到物理地址之后：向下一级回填整行，CPU 只是停顿或乱序等待。页错误是控制流改道，cache miss 只是数据慢。
\answercheck{必须保留的要点}{转换错误与数据回填是两条独立路径。}
\end{chapteranswer}

\begin{chapteranswer}{块设备队列}
CPU 把描述符写进提交队列，写 doorbell 通知设备；设备 DMA 取走描述符、搬完数据后写完成队列，再（可选）发中断。doorbell 只表示“有新请求”，不代表 I/O 完成，完成以完成队列项为准。
\answercheck{必须保留的要点}{submission 与 completion 分离，doorbell 不等于完成。}
\end{chapteranswer}

\begin{chapteranswer}{DRAM 时序计算}
一次需要开行的随机读约为 tRCD（开行）加 tCL（列出数）加突发传输，常见 DDR3/DDR4 器件合计约 30--40 ns 量级；行命中可省去 tRCD，行冲突还要先 precharge。
\answercheck{必须保留的要点}{时延由行、列、突发三段组成，不是固定常数。}
\end{chapteranswer}

\begin{chapteranswer}{UART 帧分析}
先找起始位下降沿，再按位时间中点逐位采样：1 起始位、8 数据位（低位在前）、1 停止位；还原字节后核对实测位时间与标称波特率的偏差，双方每一侧误差一般不超过约 2\%，收发合计留有约 5\% 的理论上限。
\answercheck{必须保留的要点}{位边界、采样点和误差都要核对。}
\end{chapteranswer}

\begin{chapteranswer}{组相联查找}
index 选中整个组，组内两路 tag 并行与地址 tag 比较，任一路相等且 V=1 即 hit；两个热地址各占一路后不再互相驱逐；miss 时按替换位选一路驱逐。
\answercheck{必须保留的要点}{并行比较、mux 和替换位缺一不可。}
\end{chapteranswer}

\begin{chapteranswer}{描述符环顺序}
CPU 只推进 SW 指针，设备只推进 HW 指针，末尾回绕；doorbell 只提示有新描述符，不代表完成；完成以设备写回状态位或 completion 项为准，回收前确认写回可见。
\answercheck{必须保留的要点}{指针分离、doorbell 非完成、写回可见性。}
\end{chapteranswer}

\begin{chapteranswer}{bank 交错拍数}
第 1..4 拍依次启动第 1..4 行（各占一个 bank），第 $i$ 行在第 $i$ 拍启动、第 $i+3$ 拍完成；第 5 行在第 5 拍回到 bank 1 时前一次访问刚结束，无需等待。第 8 行在第 11 拍完成，共 11 拍。单 bank 串行需 $8\times4=32$ 拍。
\answercheck{必须保留的要点}{加速来自启动流水化，单 bank 本身没有变快。}
\end{chapteranswer}

\begin{chapteranswer}{替换策略}
2 路 LRU：A miss 得 \{A,--\}；B miss 得 \{A,B\}；A hit（B 成为最久未用）；C miss 驱逐 B 得 \{A,C\}；A hit；B miss 驱逐 C 得 \{A,B\}，共 3 次 miss。直接映射下同 index 只有一个槽，每次访问都驱逐上一行，6 次全 miss。
\answercheck{必须保留的要点}{LRU 驱逐“最久未用”，不是“最早进入”。}
\end{chapteranswer}

\begin{chapteranswer}{预取演算}
无预取：每行 miss 30 拍加处理 8 拍，$16\times38=608$ 拍。next-line 预取：第 1 行第 30 拍就绪，处理第 $i$ 行的 8 拍只能隐藏预取的一部分，第 $i$ 行数据在第 $30i$ 拍就绪，稳态每行 30 拍，总 $30\times16+8=488$ 拍。处理节拍 $T\geq30$ 时预取被完全隐藏，总拍数降为 $30+16T$。
\answercheck{必须保留的要点}{预取隐藏延迟但不消除首次 miss；处理太慢时仍受 miss 间隔限制。}
\end{chapteranswer}

\begin{chapteranswer}{I2C 事务序列}
START $\rightarrow$ \code{0xA0}（\code{0x50} 左移一位、读写位 0 表示写）$\rightarrow$ 从机 ACK $\rightarrow$ \code{0x10}（寄存器号）$\rightarrow$ ACK $\rightarrow$ repeated START $\rightarrow$ \code{0xA1}（读写位 1 表示读）$\rightarrow$ ACK $\rightarrow$ 读第 1 字节 $\rightarrow$ 主机 ACK $\rightarrow$ 读第 2 字节 $\rightarrow$ 主机 NACK $\rightarrow$ STOP。
\answercheck{必须保留的要点}{先写寄存器地址，repeated START 转读；最后字节 NACK 通知从机结束。}
\end{chapteranswer}

\begin{chapteranswer}{仲裁对照}
集中式：每主设备一对请求/授权线，仲裁器并行比较，最快但线数为 $2N$ 且有单点；菊花链：授权逐级传递，线数与设备数无关，但优先级由位置固定、传播慢；轮询计数：$\lceil\log_2 N\rceil$ 根扫描线，从上次授权点继续则轮转公平，但平均扫半表最慢。片上互连把仲裁做进每个交换节点的端口。
\answercheck{必须保留的要点}{三者在线数、延迟、公平性之间取舍；仲裁没有消失。}
\end{chapteranswer}

\begin{chapteranswer}{MESI 走查}
步骤 1：A 由 I 变 S，B 保持 I（读入共享副本）；步骤 2：B 由 I 变 S，A 保持 S；步骤 3：A 广播失效，B 由 S 变 I，A 由 S 变 M，内存过期；步骤 4：B 读 miss，A 写回或转发后由 M 降为 S，B 得 S，内存恢复最新。
\answercheck{必须保留的要点}{每步最新值唯一：S 期间在内存，M 期间在 A 的 cache。}
\end{chapteranswer}
